%%%%%%%%%%%%%%%%%%%%%%%%%%%%%%%%%%%%%%%%%%%%%%%%%%%%%%%%%%%%
% 8.12 MULTIVARIATE STATISTICS
% The Composition of Advanced Mathematics
%%%%%%%%%%%%%%%%%%%%%%%%%%%%%%%%%%%%%%%%%%%%%%%%%%%%%%%%%%%%

\section{Multivariate Statistics}
\label{sec:multivariate-statistics}

\prerequisite{
Linear algebra, matrices, eigenvalues and eigenvectors, covariance,
correlation, multivariate probability, normal distributions, statistical
estimation, confidence intervals, hypothesis testing, regression, and
analysis of variance.
}

\learningobjectives{
By the end of this section, the reader should be able to:
\begin{itemize}
    \item represent multivariate observations as random vectors;
    \item define population mean vectors and covariance matrices;
    \item compute sample mean vectors and sample covariance matrices;
    \item interpret covariance and correlation matrices;
    \item understand positive semidefiniteness of covariance matrices;
    \item define the multivariate normal distribution;
    \item understand marginal and conditional normal distributions;
    \item compute Mahalanobis distance;
    \item distinguish Euclidean and covariance-adjusted distance;
    \item understand quadratic forms in multivariate statistics;
    \item perform one-sample Hotelling \(T^2\) tests;
    \item understand confidence regions for mean vectors;
    \item perform two-sample Hotelling \(T^2\) tests conceptually;
    \item understand simultaneous inference for several means;
    \item understand the purpose of MANOVA;
    \item recognize Wilks' lambda, Pillai's trace, Hotelling--Lawley
          trace, and Roy's largest root;
    \item understand principal component analysis;
    \item derive principal components from covariance eigenvectors;
    \item calculate explained-variance proportions;
    \item distinguish covariance-based and correlation-based PCA;
    \item understand dimension reduction;
    \item understand factor analysis conceptually;
    \item distinguish PCA from factor analysis;
    \item understand linear discriminant analysis;
    \item understand quadratic discriminant analysis;
    \item understand canonical correlation analysis;
    \item recognize multivariate outliers;
    \item understand covariance estimation problems in high dimensions;
    \item recognize regularized covariance estimation;
    \item understand standardization and scaling issues;
    \item understand missing-data complications in multivariate analysis;
    \item connect multivariate statistics with regression, ANOVA,
          statistical learning, and computational statistics.
\end{itemize}
}

%%%%%%%%%%%%%%%%%%%%%%%%%%%%%%%%%%%%%%%%%%%%%%%%%%%%%%%%%%%%
\subsection{Why Multivariate Statistics?}
\label{subsec:why-multivariate-statistics}
%%%%%%%%%%%%%%%%%%%%%%%%%%%%%%%%%%%%%%%%%%%%%%%%%%%%%%%%%%%%

Many datasets contain several variables measured on each observational
unit.

For example, one individual might have measurements

\[
\boxed{
\text{height},
\quad
\text{weight},
\quad
\text{blood pressure},
\quad
\text{age},
\quad
\text{cholesterol}.
}
\]

Analyzing each variable separately can miss relationships among them.

Multivariate statistics studies several variables jointly.

%%%%%%%%%%%%%%%%%%%%%%%%%%%%%%%%%%%%%%%%%%%%%%%%%%%%%%%%%%%%
\subsection{Multivariate Observation}
\label{subsec:multivariate-observation}
%%%%%%%%%%%%%%%%%%%%%%%%%%%%%%%%%%%%%%%%%%%%%%%%%%%%%%%%%%%%

A multivariate observation can be written as a vector

\[
\boxed{
\mathbf{X}
=
\begin{pmatrix}
X_1\\
X_2\\
\vdots\\
X_p
\end{pmatrix}.
}
\]

The dimension

\[
p
\]

is the number of variables measured on each unit.

A sample of size \(n\) consists of vectors

\[
\boxed{
\mathbf{X}_1,
\ldots,
\mathbf{X}_n.
}
\]

%%%%%%%%%%%%%%%%%%%%%%%%%%%%%%%%%%%%%%%%%%%%%%%%%%%%%%%%%%%%
\subsection{Data Matrix}
\label{subsec:multivariate-data-matrix}
%%%%%%%%%%%%%%%%%%%%%%%%%%%%%%%%%%%%%%%%%%%%%%%%%%%%%%%%%%%%

A multivariate dataset can be represented as

\[
\boxed{
X
=
\begin{pmatrix}
x_{11} & x_{12} & \cdots & x_{1p}\\
x_{21} & x_{22} & \cdots & x_{2p}\\
\vdots & \vdots & \ddots & \vdots\\
x_{n1} & x_{n2} & \cdots & x_{np}
\end{pmatrix}.
}
\]

Rows correspond to observational units.

Columns correspond to variables.

Thus

\[
X
\]

is an

\[
\boxed{
n\times p
}
\]

matrix.

%%%%%%%%%%%%%%%%%%%%%%%%%%%%%%%%%%%%%%%%%%%%%%%%%%%%%%%%%%%%
\subsection{Mean Vector}
\label{subsec:population-mean-vector}
%%%%%%%%%%%%%%%%%%%%%%%%%%%%%%%%%%%%%%%%%%%%%%%%%%%%%%%%%%%%

The population mean vector is

\[
\boxed{
\boldsymbol{\mu}
=
\mathbb{E}[\mathbf{X}]
=
\begin{pmatrix}
\mu_1\\
\mu_2\\
\vdots\\
\mu_p
\end{pmatrix}.
}
\]

Each component is

\[
\boxed{
\mu_j
=
\mathbb{E}[X_j].
}
\]

Thus the mean vector records the mean of every variable.

%%%%%%%%%%%%%%%%%%%%%%%%%%%%%%%%%%%%%%%%%%%%%%%%%%%%%%%%%%%%
\subsection{Covariance Matrix}
\label{subsec:population-covariance-matrix}
%%%%%%%%%%%%%%%%%%%%%%%%%%%%%%%%%%%%%%%%%%%%%%%%%%%%%%%%%%%%

The covariance matrix is

\[
\boxed{
\Sigma
=
\operatorname{Cov}(\mathbf{X})
=
\mathbb{E}
\left[
(\mathbf{X}-\boldsymbol{\mu})
(\mathbf{X}-\boldsymbol{\mu})^T
\right].
}
\]

It has the form

\[
\boxed{
\Sigma
=
\begin{pmatrix}
\sigma_{11} & \sigma_{12} & \cdots & \sigma_{1p}\\
\sigma_{21} & \sigma_{22} & \cdots & \sigma_{2p}\\
\vdots & \vdots & \ddots & \vdots\\
\sigma_{p1} & \sigma_{p2} & \cdots & \sigma_{pp}
\end{pmatrix}.
}
\]

%%%%%%%%%%%%%%%%%%%%%%%%%%%%%%%%%%%%%%%%%%%%%%%%%%%%%%%%%%%%
\subsection{Diagonal Entries of the Covariance Matrix}
\label{subsec:covariance-matrix-diagonal}
%%%%%%%%%%%%%%%%%%%%%%%%%%%%%%%%%%%%%%%%%%%%%%%%%%%%%%%%%%%%

The diagonal entries are variances:

\[
\boxed{
\sigma_{jj}
=
\operatorname{Var}(X_j).
}
\]

Thus

\[
\Sigma
\]

contains all individual variable variances on its diagonal.

%%%%%%%%%%%%%%%%%%%%%%%%%%%%%%%%%%%%%%%%%%%%%%%%%%%%%%%%%%%%
\subsection{Off-Diagonal Entries}
\label{subsec:covariance-matrix-offdiagonal}
%%%%%%%%%%%%%%%%%%%%%%%%%%%%%%%%%%%%%%%%%%%%%%%%%%%%%%%%%%%%

For

\[
j\neq k,
\]

the entry

\[
\boxed{
\sigma_{jk}
=
\operatorname{Cov}(X_j,X_k)
}
\]

measures linear association between variables \(X_j\) and \(X_k\).

Because covariance is symmetric,

\[
\boxed{
\sigma_{jk}
=
\sigma_{kj}.
}
\]

Therefore,

\[
\boxed{
\Sigma^T=\Sigma.
}
\]

%%%%%%%%%%%%%%%%%%%%%%%%%%%%%%%%%%%%%%%%%%%%%%%%%%%%%%%%%%%%
\subsection{Positive Semidefinite Covariance Matrices}
\label{subsec:covariance-positive-semidefinite}
%%%%%%%%%%%%%%%%%%%%%%%%%%%%%%%%%%%%%%%%%%%%%%%%%%%%%%%%%%%%

For any vector

\[
\mathbf{a}\in\mathbb{R}^p,
\]

consider the linear combination

\[
\mathbf{a}^T\mathbf{X}.
\]

Its variance is

\[
\boxed{
\operatorname{Var}
(
\mathbf{a}^T\mathbf{X}
)
=
\mathbf{a}^T
\Sigma
\mathbf{a}.
}
\]

Since variance cannot be negative,

\[
\boxed{
\mathbf{a}^T
\Sigma
\mathbf{a}
\geq0.
}
\]

Thus every covariance matrix is positive semidefinite.

%%%%%%%%%%%%%%%%%%%%%%%%%%%%%%%%%%%%%%%%%%%%%%%%%%%%%%%%%%%%
\subsection{Positive Definite Covariance Matrix}
\label{subsec:positive-definite-covariance}
%%%%%%%%%%%%%%%%%%%%%%%%%%%%%%%%%%%%%%%%%%%%%%%%%%%%%%%%%%%%

If

\[
\mathbf{a}^T
\Sigma
\mathbf{a}>0
\]

for every nonzero \(\mathbf{a}\), then \(\Sigma\) is positive definite.

In that case,

\[
\boxed{
\det(\Sigma)>0
}
\]

and

\[
\boxed{
\Sigma^{-1}
}
\]

exists.

Singularity means some linear combination of the variables has zero
variance.

%%%%%%%%%%%%%%%%%%%%%%%%%%%%%%%%%%%%%%%%%%%%%%%%%%%%%%%%%%%%
\subsection{Correlation Matrix}
\label{subsec:multivariate-correlation-matrix}
%%%%%%%%%%%%%%%%%%%%%%%%%%%%%%%%%%%%%%%%%%%%%%%%%%%%%%%%%%%%

The correlation between variables \(X_j\) and \(X_k\) is

\[
\boxed{
\rho_{jk}
=
\frac{
\sigma_{jk}
}{
\sqrt{
\sigma_{jj}\sigma_{kk}
}
}.
}
\]

The correlation matrix is

\[
\boxed{
R
=
(\rho_{jk}).
}
\]

Its diagonal entries satisfy

\[
\boxed{
\rho_{jj}=1.
}
\]

%%%%%%%%%%%%%%%%%%%%%%%%%%%%%%%%%%%%%%%%%%%%%%%%%%%%%%%%%%%%
\subsection{Covariance Versus Correlation Matrices}
\label{subsec:covariance-versus-correlation-matrix}
%%%%%%%%%%%%%%%%%%%%%%%%%%%%%%%%%%%%%%%%%%%%%%%%%%%%%%%%%%%%

Covariance depends on measurement units.

Correlation is dimensionless.

Therefore covariance-based analysis is sensitive to scale, whereas
correlation-based analysis effectively standardizes the variables.

This distinction becomes especially important in principal component
analysis.

%%%%%%%%%%%%%%%%%%%%%%%%%%%%%%%%%%%%%%%%%%%%%%%%%%%%%%%%%%%%
\subsection{Sample Mean Vector}
\label{subsec:sample-mean-vector}
%%%%%%%%%%%%%%%%%%%%%%%%%%%%%%%%%%%%%%%%%%%%%%%%%%%%%%%%%%%%

Given observations

\[
\mathbf{X}_1,\ldots,\mathbf{X}_n,
\]

the sample mean vector is

\[
\boxed{
\overline{\mathbf{X}}
=
\frac1n
\sum_{i=1}^{n}
\mathbf{X}_i.
}
\]

Componentwise,

\[
\boxed{
\overline{X}_j
=
\frac1n
\sum_{i=1}^{n}
X_{ij}.
}
\]

%%%%%%%%%%%%%%%%%%%%%%%%%%%%%%%%%%%%%%%%%%%%%%%%%%%%%%%%%%%%
\subsection{Sample Covariance Matrix}
\label{subsec:sample-covariance-matrix}
%%%%%%%%%%%%%%%%%%%%%%%%%%%%%%%%%%%%%%%%%%%%%%%%%%%%%%%%%%%%

The sample covariance matrix is

\[
\boxed{
S
=
\frac1{n-1}
\sum_{i=1}^{n}
(
\mathbf{X}_i-\overline{\mathbf{X}}
)
(
\mathbf{X}_i-\overline{\mathbf{X}}
)^T.
}
\]

The entry in row \(j\), column \(k\) is

\[
\boxed{
s_{jk}
=
\frac1{n-1}
\sum_{i=1}^{n}
(X_{ij}-\overline{X}_j)
(X_{ik}-\overline{X}_k).
}
\]

%%%%%%%%%%%%%%%%%%%%%%%%%%%%%%%%%%%%%%%%%%%%%%%%%%%%%%%%%%%%
\subsection{Sample Correlation Matrix}
\label{subsec:sample-correlation-matrix}
%%%%%%%%%%%%%%%%%%%%%%%%%%%%%%%%%%%%%%%%%%%%%%%%%%%%%%%%%%%%

The sample correlation is

\[
\boxed{
r_{jk}
=
\frac{
s_{jk}
}{
\sqrt{
s_{jj}s_{kk}
}
}.
}
\]

Collecting all pairwise sample correlations gives the sample correlation
matrix

\[
\boxed{
R=(r_{jk}).
}
\]

%%%%%%%%%%%%%%%%%%%%%%%%%%%%%%%%%%%%%%%%%%%%%%%%%%%%%%%%%%%%
\subsection{Worked Example: Sample Mean Vector}
\label{subsec:worked-sample-mean-vector}
%%%%%%%%%%%%%%%%%%%%%%%%%%%%%%%%%%%%%%%%%%%%%%%%%%%%%%%%%%%%

\begin{workedexamplebox}{Two Variables}

Suppose the observations are

\[
\mathbf{x}_1=
\begin{pmatrix}
1\\
4
\end{pmatrix},
\qquad
\mathbf{x}_2=
\begin{pmatrix}
3\\
6
\end{pmatrix},
\qquad
\mathbf{x}_3=
\begin{pmatrix}
5\\
8
\end{pmatrix}.
\]

Then

\[
\overline{\mathbf{x}}
=
\frac13
\left[
\begin{pmatrix}
1\\
4
\end{pmatrix}
+
\begin{pmatrix}
3\\
6
\end{pmatrix}
+
\begin{pmatrix}
5\\
8
\end{pmatrix}
\right].
\]

Therefore,

\begin{answerbox}
\[
\boxed{
\overline{\mathbf{x}}
=
\begin{pmatrix}
3\\
6
\end{pmatrix}.
}
\]
\end{answerbox}

\end{workedexamplebox}

%%%%%%%%%%%%%%%%%%%%%%%%%%%%%%%%%%%%%%%%%%%%%%%%%%%%%%%%%%%%
\subsection{Linear Transformations of Random Vectors}
\label{subsec:linear-transformations-random-vectors}
%%%%%%%%%%%%%%%%%%%%%%%%%%%%%%%%%%%%%%%%%%%%%%%%%%%%%%%%%%%%

Let

\[
\mathbf{Y}
=
A\mathbf{X}+\mathbf{b}.
\]

Then

\[
\boxed{
\mathbb{E}[\mathbf{Y}]
=
A\boldsymbol{\mu}+\mathbf{b}.
}
\]

The covariance matrix transforms as

\[
\boxed{
\operatorname{Cov}(\mathbf{Y})
=
A\Sigma A^T.
}
\]

This identity is central throughout multivariate statistics.

%%%%%%%%%%%%%%%%%%%%%%%%%%%%%%%%%%%%%%%%%%%%%%%%%%%%%%%%%%%%
\subsection{Variance of a Linear Combination}
\label{subsec:variance-linear-combination-multivariate}
%%%%%%%%%%%%%%%%%%%%%%%%%%%%%%%%%%%%%%%%%%%%%%%%%%%%%%%%%%%%

For scalar

\[
Y
=
\mathbf{a}^T\mathbf{X},
\]

we have

\[
\boxed{
\mathbb{E}[Y]
=
\mathbf{a}^T\boldsymbol{\mu},
}
\]

and

\[
\boxed{
\operatorname{Var}(Y)
=
\mathbf{a}^T\Sigma\mathbf{a}.
}
\]

Many multivariate procedures search for useful linear combinations of
variables.

%%%%%%%%%%%%%%%%%%%%%%%%%%%%%%%%%%%%%%%%%%%%%%%%%%%%%%%%%%%%
\subsection{Multivariate Normal Distribution}
\label{subsec:multivariate-normal-distribution}
%%%%%%%%%%%%%%%%%%%%%%%%%%%%%%%%%%%%%%%%%%%%%%%%%%%%%%%%%%%%

A random vector

\[
\mathbf{X}\in\mathbb{R}^p
\]

has a multivariate normal distribution with mean
\(\boldsymbol{\mu}\) and covariance matrix \(\Sigma\) if

\[
\boxed{
\mathbf{X}
\sim
N_p
(
\boldsymbol{\mu},
\Sigma
).
}
\]

One characterization is that every linear combination

\[
\boxed{
\mathbf{a}^T\mathbf{X}
}
\]

is normally distributed for every fixed vector \(\mathbf{a}\).

%%%%%%%%%%%%%%%%%%%%%%%%%%%%%%%%%%%%%%%%%%%%%%%%%%%%%%%%%%%%
\subsection{Multivariate Normal Density}
\label{subsec:multivariate-normal-density}
%%%%%%%%%%%%%%%%%%%%%%%%%%%%%%%%%%%%%%%%%%%%%%%%%%%%%%%%%%%%

If \(\Sigma\) is positive definite, the density is

\[
\boxed{
f(\mathbf{x})
=
\frac{
1
}{
(2\pi)^{p/2}
|\Sigma|^{1/2}
}
\exp
\left[
-\frac12
(
\mathbf{x}-\boldsymbol{\mu}
)^T
\Sigma^{-1}
(
\mathbf{x}-\boldsymbol{\mu}
)
\right].
}
\]

Here,

\[
|\Sigma|
\]

denotes the determinant of \(\Sigma\).

%%%%%%%%%%%%%%%%%%%%%%%%%%%%%%%%%%%%%%%%%%%%%%%%%%%%%%%%%%%%
\subsection{Geometry of the Multivariate Normal}
\label{subsec:geometry-multivariate-normal}
%%%%%%%%%%%%%%%%%%%%%%%%%%%%%%%%%%%%%%%%%%%%%%%%%%%%%%%%%%%%

Contours of equal multivariate normal density satisfy

\[
\boxed{
(
\mathbf{x}-\boldsymbol{\mu}
)^T
\Sigma^{-1}
(
\mathbf{x}-\boldsymbol{\mu}
)
=
c.
}
\]

These contours are ellipsoids.

The eigenvectors of \(\Sigma\) determine the ellipsoid directions.

The eigenvalues determine the squared scale along those directions.

%%%%%%%%%%%%%%%%%%%%%%%%%%%%%%%%%%%%%%%%%%%%%%%%%%%%%%%%%%%%
\subsection{Marginal Normal Distributions}
\label{subsec:marginal-multivariate-normal}
%%%%%%%%%%%%%%%%%%%%%%%%%%%%%%%%%%%%%%%%%%%%%%%%%%%%%%%%%%%%

Every subvector of a multivariate normal vector is itself multivariate
normal.

For example, partition

\[
\mathbf{X}
=
\begin{pmatrix}
\mathbf{X}_1\\
\mathbf{X}_2
\end{pmatrix},
\]

with

\[
\boldsymbol{\mu}
=
\begin{pmatrix}
\boldsymbol{\mu}_1\\
\boldsymbol{\mu}_2
\end{pmatrix},
\]

and

\[
\Sigma
=
\begin{pmatrix}
\Sigma_{11}
&
\Sigma_{12}
\\
\Sigma_{21}
&
\Sigma_{22}
\end{pmatrix}.
\]

Then

\[
\boxed{
\mathbf{X}_1
\sim
N
(
\boldsymbol{\mu}_1,
\Sigma_{11}
).
}
\]

%%%%%%%%%%%%%%%%%%%%%%%%%%%%%%%%%%%%%%%%%%%%%%%%%%%%%%%%%%%%
\subsection{Conditional Multivariate Normal Distribution}
\label{subsec:conditional-multivariate-normal}
%%%%%%%%%%%%%%%%%%%%%%%%%%%%%%%%%%%%%%%%%%%%%%%%%%%%%%%%%%%%

Under the same partition,

\[
\boxed{
\mathbf{X}_1
\mid
\mathbf{X}_2=\mathbf{x}_2
}
\]

is multivariate normal with conditional mean

\[
\boxed{
\boldsymbol{\mu}_1
+
\Sigma_{12}
\Sigma_{22}^{-1}
(
\mathbf{x}_2-\boldsymbol{\mu}_2
)
}
\]

and conditional covariance

\[
\boxed{
\Sigma_{11}
-
\Sigma_{12}
\Sigma_{22}^{-1}
\Sigma_{21}.
}
\]

This formula connects multivariate normal theory directly with linear
regression.

%%%%%%%%%%%%%%%%%%%%%%%%%%%%%%%%%%%%%%%%%%%%%%%%%%%%%%%%%%%%
\subsection{Uncorrelated Normal Variables Are Independent}
\label{subsec:uncorrelated-normal-independent}
%%%%%%%%%%%%%%%%%%%%%%%%%%%%%%%%%%%%%%%%%%%%%%%%%%%%%%%%%%%%

For a jointly multivariate normal vector, zero covariance implies
independence.

Thus if

\[
\boxed{
\operatorname{Cov}(X_j,X_k)=0
}
\]

for jointly normal components, then

\[
\boxed{
X_j
\perp X_k.
}
\]

This implication is special to particular distributional families and is
not true in general.

%%%%%%%%%%%%%%%%%%%%%%%%%%%%%%%%%%%%%%%%%%%%%%%%%%%%%%%%%%%%
\subsection{Mahalanobis Distance}
\label{subsec:mahalanobis-distance}
%%%%%%%%%%%%%%%%%%%%%%%%%%%%%%%%%%%%%%%%%%%%%%%%%%%%%%%%%%%%

The squared Mahalanobis distance from \(\mathbf{x}\) to mean
\(\boldsymbol{\mu}\) is

\[
\boxed{
D_M^2
=
(
\mathbf{x}-\boldsymbol{\mu}
)^T
\Sigma^{-1}
(
\mathbf{x}-\boldsymbol{\mu}
).
}
\]

The Mahalanobis distance is

\[
\boxed{
D_M
=
\sqrt{
D_M^2
}.
}
\]

%%%%%%%%%%%%%%%%%%%%%%%%%%%%%%%%%%%%%%%%%%%%%%%%%%%%%%%%%%%%
\subsection{Euclidean Versus Mahalanobis Distance}
\label{subsec:euclidean-versus-mahalanobis}
%%%%%%%%%%%%%%%%%%%%%%%%%%%%%%%%%%%%%%%%%%%%%%%%%%%%%%%%%%%%

Euclidean distance is

\[
\boxed{
\|\mathbf{x}-\boldsymbol{\mu}\|_2.
}
\]

It treats all coordinate directions equally.

Mahalanobis distance adjusts for:

\[
\boxed{
\text{different variable scales}
}
\]

and

\[
\boxed{
\text{correlations among variables}.
}
\]

Thus a direction with high natural variability contributes less than a
direction with low variability.

%%%%%%%%%%%%%%%%%%%%%%%%%%%%%%%%%%%%%%%%%%%%%%%%%%%%%%%%%%%%
\subsection{Worked Example: Diagonal Mahalanobis Distance}
\label{subsec:worked-mahalanobis-distance}
%%%%%%%%%%%%%%%%%%%%%%%%%%%%%%%%%%%%%%%%%%%%%%%%%%%%%%%%%%%%

\begin{workedexamplebox}{Independent Components}

Suppose

\[
\boldsymbol{\mu}
=
\begin{pmatrix}
0\\
0
\end{pmatrix},
\qquad
\Sigma
=
\begin{pmatrix}
4&0\\
0&1
\end{pmatrix},
\]

and

\[
\mathbf{x}
=
\begin{pmatrix}
2\\
1
\end{pmatrix}.
\]

Then

\[
\Sigma^{-1}
=
\begin{pmatrix}
1/4&0\\
0&1
\end{pmatrix}.
\]

Therefore,

\[
D_M^2
=
\begin{pmatrix}
2&1
\end{pmatrix}
\begin{pmatrix}
1/4&0\\
0&1
\end{pmatrix}
\begin{pmatrix}
2\\
1
\end{pmatrix}.
\]

Hence,

\begin{answerbox}
\[
\boxed{
D_M^2
=
2,
\qquad
D_M=\sqrt2.
}
\]
\end{answerbox}

\end{workedexamplebox}

%%%%%%%%%%%%%%%%%%%%%%%%%%%%%%%%%%%%%%%%%%%%%%%%%%%%%%%%%%%%
\subsection{Mahalanobis Distance Under Normality}
\label{subsec:mahalanobis-normality}
%%%%%%%%%%%%%%%%%%%%%%%%%%%%%%%%%%%%%%%%%%%%%%%%%%%%%%%%%%%%

If

\[
\mathbf{X}
\sim
N_p
(
\boldsymbol{\mu},
\Sigma
),
\]

then

\[
\boxed{
(
\mathbf{X}-\boldsymbol{\mu}
)^T
\Sigma^{-1}
(
\mathbf{X}-\boldsymbol{\mu}
)
\sim
\chi^2_p.
}
\]

This makes Mahalanobis distance useful for identifying unusually distant
multivariate observations.

%%%%%%%%%%%%%%%%%%%%%%%%%%%%%%%%%%%%%%%%%%%%%%%%%%%%%%%%%%%%
\subsection{Multivariate Outliers}
\label{subsec:multivariate-outliers}
%%%%%%%%%%%%%%%%%%%%%%%%%%%%%%%%%%%%%%%%%%%%%%%%%%%%%%%%%%%%

An observation may appear ordinary on every variable separately but
still represent an unusual multivariate combination.

For example, two measurements may individually be common but rarely occur
together.

Mahalanobis distance can detect this joint unusualness.

%%%%%%%%%%%%%%%%%%%%%%%%%%%%%%%%%%%%%%%%%%%%%%%%%%%%%%%%%%%%
\subsection{One-Sample Mean-Vector Testing}
\label{subsec:one-sample-mean-vector-testing}
%%%%%%%%%%%%%%%%%%%%%%%%%%%%%%%%%%%%%%%%%%%%%%%%%%%%%%%%%%%%

Suppose

\[
\mathbf{X}_1,\ldots,\mathbf{X}_n
\]

are i.i.d. observations from a \(p\)-dimensional population.

Consider

\[
\boxed{
H_0:
\boldsymbol{\mu}
=
\boldsymbol{\mu}_0.
}
\]

Testing each component separately creates a multiple-testing problem.

Hotelling's \(T^2\) provides a joint test.

%%%%%%%%%%%%%%%%%%%%%%%%%%%%%%%%%%%%%%%%%%%%%%%%%%%%%%%%%%%%
\subsection{Hotelling's \(T^2\) Statistic}
\label{subsec:hotelling-t-squared}
%%%%%%%%%%%%%%%%%%%%%%%%%%%%%%%%%%%%%%%%%%%%%%%%%%%%%%%%%%%%

The one-sample statistic is

\[
\boxed{
T^2
=
n
(
\overline{\mathbf{X}}
-
\boldsymbol{\mu}_0
)^T
S^{-1}
(
\overline{\mathbf{X}}
-
\boldsymbol{\mu}_0
).
}
\]

This is a multivariate analogue of the one-sample squared \(t\)-statistic.

%%%%%%%%%%%%%%%%%%%%%%%%%%%%%%%%%%%%%%%%%%%%%%%%%%%%%%%%%%%%
\subsection{Hotelling \(T^2\) Reference Distribution}
\label{subsec:hotelling-reference-distribution}
%%%%%%%%%%%%%%%%%%%%%%%%%%%%%%%%%%%%%%%%%%%%%%%%%%%%%%%%%%%%

Under multivariate normal sampling and

\[
H_0:
\boldsymbol{\mu}
=
\boldsymbol{\mu}_0,
\]

\[
\boxed{
\frac{
n-p
}{
p(n-1)
}
T^2
\sim
F_{p,n-p},
}
\]

provided

\[
\boxed{
n>p.
}
\]

Thus the test uses an \(F\)-distribution after rescaling.

%%%%%%%%%%%%%%%%%%%%%%%%%%%%%%%%%%%%%%%%%%%%%%%%%%%%%%%%%%%%
\subsection{Reduction to the One-Sample \(t\)-Test}
\label{subsec:hotelling-reduction-t}
%%%%%%%%%%%%%%%%%%%%%%%%%%%%%%%%%%%%%%%%%%%%%%%%%%%%%%%%%%%%

When

\[
p=1,
\]

Hotelling's statistic becomes

\[
T^2
=
n
\frac{
(\overline{X}-\mu_0)^2
}{
S^2
}.
\]

But

\[
t
=
\frac{
\overline{X}-\mu_0
}{
S/\sqrt{n}
}.
\]

Therefore,

\[
\boxed{
T^2=t^2.
}
\]

Thus Hotelling's test generalizes the ordinary one-sample \(t\)-test.

%%%%%%%%%%%%%%%%%%%%%%%%%%%%%%%%%%%%%%%%%%%%%%%%%%%%%%%%%%%%
\subsection{Worked Example: Hotelling Statistic Setup}
\label{subsec:worked-hotelling-setup}
%%%%%%%%%%%%%%%%%%%%%%%%%%%%%%%%%%%%%%%%%%%%%%%%%%%%%%%%%%%%

\begin{workedexamplebox}{Testing a Two-Dimensional Mean}

Suppose

\[
n=20,
\]

\[
\overline{\mathbf{x}}
=
\begin{pmatrix}
5\\
8
\end{pmatrix},
\]

and the null mean is

\[
\boldsymbol{\mu}_0
=
\begin{pmatrix}
4\\
6
\end{pmatrix}.
\]

Suppose

\[
S^{-1}
=
\begin{pmatrix}
0.5&-0.1\\
-0.1&0.4
\end{pmatrix}.
\]

Then

\[
\overline{\mathbf{x}}
-
\boldsymbol{\mu}_0
=
\begin{pmatrix}
1\\
2
\end{pmatrix}.
\]

Therefore,

\begin{answerbox}
\[
\boxed{
T^2
=
20
\begin{pmatrix}
1&2
\end{pmatrix}
\begin{pmatrix}
0.5&-0.1\\
-0.1&0.4
\end{pmatrix}
\begin{pmatrix}
1\\
2
\end{pmatrix}.
}
\]
\end{answerbox}

\end{workedexamplebox}

%%%%%%%%%%%%%%%%%%%%%%%%%%%%%%%%%%%%%%%%%%%%%%%%%%%%%%%%%%%%
\subsection{Confidence Region for a Mean Vector}
\label{subsec:confidence-region-mean-vector}
%%%%%%%%%%%%%%%%%%%%%%%%%%%%%%%%%%%%%%%%%%%%%%%%%%%%%%%%%%%%

A multivariate confidence set for \(\boldsymbol{\mu}\) can be written

\[
\boxed{
n
(
\overline{\mathbf{X}}
-
\boldsymbol{\mu}
)^T
S^{-1}
(
\overline{\mathbf{X}}
-
\boldsymbol{\mu}
)
\leq
c.
}
\]

Under multivariate normality,

\[
\boxed{
c
=
\frac{
p(n-1)
}{
n-p
}
F_{p,n-p;1-\alpha}.
}
\]

This produces an ellipsoidal confidence region.

%%%%%%%%%%%%%%%%%%%%%%%%%%%%%%%%%%%%%%%%%%%%%%%%%%%%%%%%%%%%
\subsection{Simultaneous Confidence Intervals}
\label{subsec:simultaneous-multivariate-ci}
%%%%%%%%%%%%%%%%%%%%%%%%%%%%%%%%%%%%%%%%%%%%%%%%%%%%%%%%%%%%

The confidence ellipsoid can be used to derive simultaneous intervals for
linear combinations

\[
\mathbf{a}^T\boldsymbol{\mu}.
\]

A general form is

\[
\boxed{
\mathbf{a}^T
\overline{\mathbf{X}}
\pm
\sqrt{
\frac{c}{n}
\mathbf{a}^TS\mathbf{a}
}.
}
\]

These intervals protect joint coverage over a specified family of linear
combinations.

%%%%%%%%%%%%%%%%%%%%%%%%%%%%%%%%%%%%%%%%%%%%%%%%%%%%%%%%%%%%
\subsection{Bonferroni Simultaneous Intervals}
\label{subsec:bonferroni-multivariate-ci}
%%%%%%%%%%%%%%%%%%%%%%%%%%%%%%%%%%%%%%%%%%%%%%%%%%%%%%%%%%%%

If interest is restricted to the \(p\) individual mean components,
Bonferroni intervals provide a simple alternative.

For component \(j\),

\[
\boxed{
\overline{X}_j
\pm
t^*
\sqrt{
\frac{s_{jj}}{n}
},
}
\]

where \(t^*\) uses a tail probability adjusted for the number of
simultaneous intervals.

Bonferroni methods can be conservative but are easy to construct.

%%%%%%%%%%%%%%%%%%%%%%%%%%%%%%%%%%%%%%%%%%%%%%%%%%%%%%%%%%%%
\subsection{Two-Sample Mean-Vector Comparison}
\label{subsec:two-sample-multivariate-mean}
%%%%%%%%%%%%%%%%%%%%%%%%%%%%%%%%%%%%%%%%%%%%%%%%%%%%%%%%%%%%

Suppose two independent multivariate populations have mean vectors

\[
\boldsymbol{\mu}_1
\]

and

\[
\boldsymbol{\mu}_2.
\]

A common hypothesis is

\[
\boxed{
H_0:
\boldsymbol{\mu}_1
=
\boldsymbol{\mu}_2.
}
\]

The point estimator is

\[
\boxed{
\overline{\mathbf{X}}_1
-
\overline{\mathbf{X}}_2.
}
\]

%%%%%%%%%%%%%%%%%%%%%%%%%%%%%%%%%%%%%%%%%%%%%%%%%%%%%%%%%%%%
\subsection{Pooled Covariance Matrix}
\label{subsec:multivariate-pooled-covariance}
%%%%%%%%%%%%%%%%%%%%%%%%%%%%%%%%%%%%%%%%%%%%%%%%%%%%%%%%%%%%

Under the assumption of common covariance

\[
\Sigma_1=\Sigma_2=\Sigma,
\]

the pooled sample covariance matrix is

\[
\boxed{
S_p
=
\frac{
(n_1-1)S_1
+
(n_2-1)S_2
}{
n_1+n_2-2
}.
}
\]

%%%%%%%%%%%%%%%%%%%%%%%%%%%%%%%%%%%%%%%%%%%%%%%%%%%%%%%%%%%%
\subsection{Two-Sample Hotelling \(T^2\)}
\label{subsec:two-sample-hotelling}
%%%%%%%%%%%%%%%%%%%%%%%%%%%%%%%%%%%%%%%%%%%%%%%%%%%%%%%%%%%%

The two-sample statistic is

\[
\boxed{
T^2
=
\frac{
n_1n_2
}{
n_1+n_2
}
(
\overline{\mathbf{X}}_1
-
\overline{\mathbf{X}}_2
)^T
S_p^{-1}
(
\overline{\mathbf{X}}_1
-
\overline{\mathbf{X}}_2
).
}
\]

Under multivariate normality and equal covariance matrices, this can be
transformed to an \(F\)-statistic.

%%%%%%%%%%%%%%%%%%%%%%%%%%%%%%%%%%%%%%%%%%%%%%%%%%%%%%%%%%%%
\subsection{Two-Sample Hotelling Reference Distribution}
\label{subsec:two-sample-hotelling-f}
%%%%%%%%%%%%%%%%%%%%%%%%%%%%%%%%%%%%%%%%%%%%%%%%%%%%%%%%%%%%

Under the classical null,

\[
\boxed{
\frac{
n_1+n_2-p-1
}{
p(n_1+n_2-2)
}
T^2
\sim
F_{
p,\,
n_1+n_2-p-1
}.
}
\]

This requires sufficient sample size for the pooled covariance matrix to
be nonsingular.

%%%%%%%%%%%%%%%%%%%%%%%%%%%%%%%%%%%%%%%%%%%%%%%%%%%%%%%%%%%%
\subsection{Why Test Mean Vectors Jointly?}
\label{subsec:why-test-mean-vectors-jointly}
%%%%%%%%%%%%%%%%%%%%%%%%%%%%%%%%%%%%%%%%%%%%%%%%%%%%%%%%%%%%

Suppose five correlated outcomes are measured.

Testing five separate hypotheses ignores their covariance structure and
creates multiplicity.

A joint multivariate test can detect a combined pattern such as:

\[
\boxed{
\text{small coordinated changes across several variables}.
}
\]

This may be scientifically important even when no individual univariate
test is extremely strong.

%%%%%%%%%%%%%%%%%%%%%%%%%%%%%%%%%%%%%%%%%%%%%%%%%%%%%%%%%%%%
\subsection{MANOVA}
\label{subsec:manova}
%%%%%%%%%%%%%%%%%%%%%%%%%%%%%%%%%%%%%%%%%%%%%%%%%%%%%%%%%%%%

Multivariate analysis of variance, abbreviated MANOVA, extends ANOVA to
multiple response variables.

Instead of a scalar response,

\[
Y,
\]

each experimental unit has a response vector

\[
\boxed{
\mathbf{Y}
=
\begin{pmatrix}
Y_1\\
\vdots\\
Y_p
\end{pmatrix}.
}
\]

The objective is to test whether groups differ jointly across the
response vector.

%%%%%%%%%%%%%%%%%%%%%%%%%%%%%%%%%%%%%%%%%%%%%%%%%%%%%%%%%%%%
\subsection{MANOVA Model}
\label{subsec:manova-model}
%%%%%%%%%%%%%%%%%%%%%%%%%%%%%%%%%%%%%%%%%%%%%%%%%%%%%%%%%%%%

A multivariate linear model can be written

\[
\boxed{
Y
=
XB+E,
}
\]

where

\[
Y
\]

is an

\[
n\times p
\]

matrix of responses,

\[
X
\]

is the design matrix,

\[
B
\]

contains regression or treatment coefficients, and

\[
E
\]

contains multivariate errors.

%%%%%%%%%%%%%%%%%%%%%%%%%%%%%%%%%%%%%%%%%%%%%%%%%%%%%%%%%%%%
\subsection{Hypothesis and Error Matrices}
\label{subsec:manova-h-e-matrices}
%%%%%%%%%%%%%%%%%%%%%%%%%%%%%%%%%%%%%%%%%%%%%%%%%%%%%%%%%%%%

MANOVA partitions multivariate sums of squares and cross-products into
matrices.

The hypothesis matrix is often denoted

\[
\boxed{
H,
}
\]

while the error matrix is

\[
\boxed{
E.
}
\]

These generalize the scalar ANOVA quantities

\[
SSR
\]

and

\[
SSE.
\]

%%%%%%%%%%%%%%%%%%%%%%%%%%%%%%%%%%%%%%%%%%%%%%%%%%%%%%%%%%%%
\subsection{Wilks' Lambda}
\label{subsec:wilks-lambda-manova}
%%%%%%%%%%%%%%%%%%%%%%%%%%%%%%%%%%%%%%%%%%%%%%%%%%%%%%%%%%%%

One MANOVA statistic is Wilks' lambda:

\[
\boxed{
\Lambda
=
\frac{
|E|
}{
|E+H|
}.
}
\]

Values near \(1\) indicate little additional variation explained by the
tested effect.

Small values provide stronger evidence against the null.

%%%%%%%%%%%%%%%%%%%%%%%%%%%%%%%%%%%%%%%%%%%%%%%%%%%%%%%%%%%%
\subsection{Pillai's Trace}
\label{subsec:pillais-trace}
%%%%%%%%%%%%%%%%%%%%%%%%%%%%%%%%%%%%%%%%%%%%%%%%%%%%%%%%%%%%

Another MANOVA statistic is Pillai's trace:

\[
\boxed{
V
=
\operatorname{tr}
\left[
H(H+E)^{-1}
\right].
}
\]

Larger values indicate stronger multivariate group effects.

Pillai's trace is often considered relatively robust among classical
MANOVA criteria.

%%%%%%%%%%%%%%%%%%%%%%%%%%%%%%%%%%%%%%%%%%%%%%%%%%%%%%%%%%%%
\subsection{Hotelling--Lawley Trace}
\label{subsec:hotelling-lawley-trace}
%%%%%%%%%%%%%%%%%%%%%%%%%%%%%%%%%%%%%%%%%%%%%%%%%%%%%%%%%%%%

The Hotelling--Lawley trace is

\[
\boxed{
T
=
\operatorname{tr}
(
E^{-1}H
).
}
\]

Larger values indicate greater separation relative to within-group
variation.

%%%%%%%%%%%%%%%%%%%%%%%%%%%%%%%%%%%%%%%%%%%%%%%%%%%%%%%%%%%%
\subsection{Roy's Largest Root}
\label{subsec:roys-largest-root}
%%%%%%%%%%%%%%%%%%%%%%%%%%%%%%%%%%%%%%%%%%%%%%%%%%%%%%%%%%%%

Roy's largest root focuses on the largest eigenvalue of

\[
\boxed{
E^{-1}H.
}
\]

It emphasizes the strongest single multivariate direction of group
separation.

Different MANOVA statistics weight the multivariate directions
differently.

%%%%%%%%%%%%%%%%%%%%%%%%%%%%%%%%%%%%%%%%%%%%%%%%%%%%%%%%%%%%
\subsection{MANOVA Versus Separate ANOVAs}
\label{subsec:manova-versus-anovas}
%%%%%%%%%%%%%%%%%%%%%%%%%%%%%%%%%%%%%%%%%%%%%%%%%%%%%%%%%%%%

Separate ANOVAs ignore relationships among response variables.

MANOVA incorporates their covariance structure.

However, a significant MANOVA does not automatically identify:

\[
\boxed{
\text{which outcomes differ},
}
\]

or

\[
\boxed{
\text{which groups differ}.
}
\]

Follow-up analyses are still required.

%%%%%%%%%%%%%%%%%%%%%%%%%%%%%%%%%%%%%%%%%%%%%%%%%%%%%%%%%%%%
\subsection{MANOVA Assumptions}
\label{subsec:manova-assumptions}
%%%%%%%%%%%%%%%%%%%%%%%%%%%%%%%%%%%%%%%%%%%%%%%%%%%%%%%%%%%%

Classical MANOVA typically assumes:

\[
\boxed{
\text{independent observations},
}
\]

\[
\boxed{
\text{multivariate normality within groups},
}
\]

and often

\[
\boxed{
\text{equal covariance matrices across groups}.
}
\]

The analysis can be sensitive to severe violations, especially with
small samples.

%%%%%%%%%%%%%%%%%%%%%%%%%%%%%%%%%%%%%%%%%%%%%%%%%%%%%%%%%%%%
\subsection{Principal Component Analysis}
\label{subsec:principal-component-analysis}
%%%%%%%%%%%%%%%%%%%%%%%%%%%%%%%%%%%%%%%%%%%%%%%%%%%%%%%%%%%%

Principal component analysis, abbreviated PCA, transforms correlated
variables into new uncorrelated linear combinations.

The first principal component is

\[
\boxed{
Z_1
=
\mathbf{a}_1^T\mathbf{X},
}
\]

where \(\mathbf{a}_1\) is chosen to maximize

\[
\operatorname{Var}(Z_1)
\]

subject to

\[
\boxed{
\mathbf{a}_1^T\mathbf{a}_1=1.
}
\]

%%%%%%%%%%%%%%%%%%%%%%%%%%%%%%%%%%%%%%%%%%%%%%%%%%%%%%%%%%%%
\subsection{First Principal Component}
\label{subsec:first-principal-component}
%%%%%%%%%%%%%%%%%%%%%%%%%%%%%%%%%%%%%%%%%%%%%%%%%%%%%%%%%%%%

Since

\[
\operatorname{Var}
(
\mathbf{a}^T\mathbf{X}
)
=
\mathbf{a}^T
\Sigma
\mathbf{a},
\]

the optimization problem is

\[
\boxed{
\max_{\mathbf{a}}
\mathbf{a}^T
\Sigma
\mathbf{a}
}
\]

subject to

\[
\boxed{
\mathbf{a}^T\mathbf{a}=1.
}
\]

Using a Lagrange multiplier gives

\[
\boxed{
\Sigma\mathbf{a}
=
\lambda\mathbf{a}.
}
\]

Thus principal-component directions are eigenvectors of the covariance
matrix.

%%%%%%%%%%%%%%%%%%%%%%%%%%%%%%%%%%%%%%%%%%%%%%%%%%%%%%%%%%%%
\subsection{PCA Eigenvalues}
\label{subsec:pca-eigenvalues}
%%%%%%%%%%%%%%%%%%%%%%%%%%%%%%%%%%%%%%%%%%%%%%%%%%%%%%%%%%%%

Let

\[
\lambda_1
\geq
\lambda_2
\geq
\cdots
\geq
\lambda_p
\geq0
\]

be the eigenvalues of \(\Sigma\).

Then the variance of principal component \(j\) is

\[
\boxed{
\operatorname{Var}(Z_j)
=
\lambda_j.
}
\]

The first component therefore captures the largest possible variance
among unit-length linear combinations.

%%%%%%%%%%%%%%%%%%%%%%%%%%%%%%%%%%%%%%%%%%%%%%%%%%%%%%%%%%%%
\subsection{Orthogonality of Principal Components}
\label{subsec:pca-orthogonality}
%%%%%%%%%%%%%%%%%%%%%%%%%%%%%%%%%%%%%%%%%%%%%%%%%%%%%%%%%%%%

For a symmetric covariance matrix, eigenvectors corresponding to distinct
eigenvalues can be chosen orthonormal.

Thus

\[
\boxed{
\mathbf{a}_i^T
\mathbf{a}_j
=
0
\qquad
(i\neq j).
}
\]

Consequently,

\[
\boxed{
\operatorname{Cov}(Z_i,Z_j)=0.
}
\]

Principal components are therefore uncorrelated.

%%%%%%%%%%%%%%%%%%%%%%%%%%%%%%%%%%%%%%%%%%%%%%%%%%%%%%%%%%%%
\subsection{Total Variance in PCA}
\label{subsec:pca-total-variance}
%%%%%%%%%%%%%%%%%%%%%%%%%%%%%%%%%%%%%%%%%%%%%%%%%%%%%%%%%%%%

The total variance of the original variables is

\[
\boxed{
\operatorname{tr}(\Sigma)
=
\sum_{j=1}^{p}\sigma_{jj}.
}
\]

Because the trace equals the sum of eigenvalues,

\[
\boxed{
\operatorname{tr}(\Sigma)
=
\sum_{j=1}^{p}\lambda_j.
}
\]

Thus PCA redistributes the same total variance across orthogonal
components.

%%%%%%%%%%%%%%%%%%%%%%%%%%%%%%%%%%%%%%%%%%%%%%%%%%%%%%%%%%%%
\subsection{Explained Variance Ratio}
\label{subsec:pca-explained-variance}
%%%%%%%%%%%%%%%%%%%%%%%%%%%%%%%%%%%%%%%%%%%%%%%%%%%%%%%%%%%%

The proportion of total variance explained by principal component \(j\)
is

\[
\boxed{
\frac{
\lambda_j
}{
\sum_{k=1}^{p}\lambda_k
}.
}
\]

The cumulative proportion explained by the first \(m\) components is

\[
\boxed{
\frac{
\sum_{j=1}^{m}\lambda_j
}{
\sum_{j=1}^{p}\lambda_j
}.
}
\]

%%%%%%%%%%%%%%%%%%%%%%%%%%%%%%%%%%%%%%%%%%%%%%%%%%%%%%%%%%%%
\subsection{Worked Example: PCA Variance}
\label{subsec:worked-pca-variance}
%%%%%%%%%%%%%%%%%%%%%%%%%%%%%%%%%%%%%%%%%%%%%%%%%%%%%%%%%%%%

\begin{workedexamplebox}{Explained Variance}

Suppose a covariance matrix has eigenvalues

\[
\lambda_1=8,
\qquad
\lambda_2=3,
\qquad
\lambda_3=1.
\]

The total variance is

\[
8+3+1=12.
\]

The first component explains

\[
\frac8{12}
=
\frac23.
\]

The first two components explain

\[
\frac{8+3}{12}
=
\frac{11}{12}.
\]

\begin{answerbox}
\[
\boxed{
PC_1:
66.7\%,
\qquad
PC_1+PC_2:
91.7\%.
}
\]
\end{answerbox}

\end{workedexamplebox}

%%%%%%%%%%%%%%%%%%%%%%%%%%%%%%%%%%%%%%%%%%%%%%%%%%%%%%%%%%%%
\subsection{PCA Scores}
\label{subsec:pca-scores}
%%%%%%%%%%%%%%%%%%%%%%%%%%%%%%%%%%%%%%%%%%%%%%%%%%%%%%%%%%%%

For centered observation

\[
\mathbf{x}_i-\overline{\mathbf{x}},
\]

the score on component \(j\) is

\[
\boxed{
z_{ij}
=
\mathbf{a}_j^T
(
\mathbf{x}_i-\overline{\mathbf{x}}
).
}
\]

These scores represent each observation in the lower-dimensional
principal-component coordinate system.

%%%%%%%%%%%%%%%%%%%%%%%%%%%%%%%%%%%%%%%%%%%%%%%%%%%%%%%%%%%%
\subsection{PCA Loadings}
\label{subsec:pca-loadings}
%%%%%%%%%%%%%%%%%%%%%%%%%%%%%%%%%%%%%%%%%%%%%%%%%%%%%%%%%%%%

The term loading is used for coefficients describing how original
variables contribute to components.

Depending on convention, loadings may refer directly to eigenvector
entries or scaled coefficients such as

\[
\boxed{
\sqrt{\lambda_j}\,a_{kj}.
}
\]

Therefore the precise software convention should be checked.

%%%%%%%%%%%%%%%%%%%%%%%%%%%%%%%%%%%%%%%%%%%%%%%%%%%%%%%%%%%%
\subsection{Covariance PCA Versus Correlation PCA}
\label{subsec:covariance-versus-correlation-pca}
%%%%%%%%%%%%%%%%%%%%%%%%%%%%%%%%%%%%%%%%%%%%%%%%%%%%%%%%%%%%

Covariance PCA uses the sample covariance matrix

\[
S.
\]

Correlation PCA uses the standardized variables and therefore the
correlation matrix

\[
R.
\]

If variables have very different units or scales, covariance PCA may be
dominated by the variables with the largest variances.

%%%%%%%%%%%%%%%%%%%%%%%%%%%%%%%%%%%%%%%%%%%%%%%%%%%%%%%%%%%%
\subsection{Standardization Before PCA}
\label{subsec:standardization-before-pca}
%%%%%%%%%%%%%%%%%%%%%%%%%%%%%%%%%%%%%%%%%%%%%%%%%%%%%%%%%%%%

A standardized variable is

\[
\boxed{
Z_j
=
\frac{
X_j-\mu_j
}{
\sigma_j
}.
}
\]

For sample data,

\[
\boxed{
z_{ij}
=
\frac{
x_{ij}-\overline{x}_j
}{
s_j
}.
}
\]

PCA on standardized data is equivalent to PCA based on the correlation
matrix.

%%%%%%%%%%%%%%%%%%%%%%%%%%%%%%%%%%%%%%%%%%%%%%%%%%%%%%%%%%%%
\subsection{PCA as Dimension Reduction}
\label{subsec:pca-dimension-reduction}
%%%%%%%%%%%%%%%%%%%%%%%%%%%%%%%%%%%%%%%%%%%%%%%%%%%%%%%%%%%%

Suppose

\[
p=20,
\]

but the first three principal components explain most of the variance.

Then the data may be approximated in a three-dimensional representation.

Thus PCA maps

\[
\boxed{
\mathbb{R}^p
\longrightarrow
\mathbb{R}^m,
\qquad
m<p.
}
\]

This is unsupervised dimension reduction.

%%%%%%%%%%%%%%%%%%%%%%%%%%%%%%%%%%%%%%%%%%%%%%%%%%%%%%%%%%%%
\subsection{PCA Reconstruction}
\label{subsec:pca-reconstruction}
%%%%%%%%%%%%%%%%%%%%%%%%%%%%%%%%%%%%%%%%%%%%%%%%%%%%%%%%%%%%

Let

\[
A_m
=
\begin{pmatrix}
\mathbf{a}_1&
\cdots&
\mathbf{a}_m
\end{pmatrix}.
\]

The lower-dimensional score vector is

\[
\boxed{
\mathbf{z}
=
A_m^T
(
\mathbf{x}-\boldsymbol{\mu}
).
}
\]

An approximation to the original observation is

\[
\boxed{
\widehat{\mathbf{x}}
=
\boldsymbol{\mu}
+
A_m\mathbf{z}.
}
\]

%%%%%%%%%%%%%%%%%%%%%%%%%%%%%%%%%%%%%%%%%%%%%%%%%%%%%%%%%%%%
\subsection{Optimal Reconstruction Property of PCA}
\label{subsec:pca-optimal-reconstruction}
%%%%%%%%%%%%%%%%%%%%%%%%%%%%%%%%%%%%%%%%%%%%%%%%%%%%%%%%%%%%

Among all \(m\)-dimensional linear subspaces, the principal-component
subspace minimizes mean squared reconstruction error under the usual
Euclidean criterion.

Thus PCA is not merely an arbitrary rotation.

It provides an optimal linear low-rank representation in this sense.

%%%%%%%%%%%%%%%%%%%%%%%%%%%%%%%%%%%%%%%%%%%%%%%%%%%%%%%%%%%%
\subsection{Scree Plot}
\label{subsec:scree-plot}
%%%%%%%%%%%%%%%%%%%%%%%%%%%%%%%%%%%%%%%%%%%%%%%%%%%%%%%%%%%%

A scree plot graphs

\[
\boxed{
\lambda_j
}
\]

against component number \(j\).

A sharp bend or ``elbow'' may suggest a natural dimension for retaining
components.

However, the elbow rule is subjective and should be combined with:

\[
\boxed{
\text{explained variance},
}
\]

\[
\boxed{
\text{scientific interpretability},
}
\]

and

\[
\boxed{
\text{predictive goals}.
}
\]

%%%%%%%%%%%%%%%%%%%%%%%%%%%%%%%%%%%%%%%%%%%%%%%%%%%%%%%%%%%%
\subsection{PCA Does Not Use the Response Variable}
\label{subsec:pca-unsupervised}
%%%%%%%%%%%%%%%%%%%%%%%%%%%%%%%%%%%%%%%%%%%%%%%%%%%%%%%%%%%%

Ordinary PCA chooses directions that explain variation in the predictors.

It does not use an outcome variable to determine which directions are
predictively useful.

Therefore a low-variance principal component can still be important for
predicting a response.

PCA is fundamentally unsupervised.

%%%%%%%%%%%%%%%%%%%%%%%%%%%%%%%%%%%%%%%%%%%%%%%%%%%%%%%%%%%%
\subsection{Factor Analysis}
\label{subsec:factor-analysis}
%%%%%%%%%%%%%%%%%%%%%%%%%%%%%%%%%%%%%%%%%%%%%%%%%%%%%%%%%%%%

Factor analysis models observed variables through a smaller number of
unobserved latent factors.

A common model is

\[
\boxed{
\mathbf{X}
=
\boldsymbol{\mu}
+
\Lambda\mathbf{F}
+
\boldsymbol{\varepsilon},
}
\]

where

\[
\mathbf{F}
\]

is a vector of latent factors and

\[
\Lambda
\]

is the loading matrix.

%%%%%%%%%%%%%%%%%%%%%%%%%%%%%%%%%%%%%%%%%%%%%%%%%%%%%%%%%%%%
\subsection{Factor Model Covariance}
\label{subsec:factor-model-covariance}
%%%%%%%%%%%%%%%%%%%%%%%%%%%%%%%%%%%%%%%%%%%%%%%%%%%%%%%%%%%%

Under a common standardized factor model,

\[
\operatorname{Cov}(\mathbf{F})=I
\]

and

\[
\operatorname{Cov}(\boldsymbol{\varepsilon})
=
\Psi,
\]

where \(\Psi\) is diagonal.

Then

\[
\boxed{
\Sigma
=
\Lambda\Lambda^T
+
\Psi.
}
\]

Thus covariance is explained through:

\[
\boxed{
\text{common factors}
+
\text{variable-specific variation}.
}
\]

%%%%%%%%%%%%%%%%%%%%%%%%%%%%%%%%%%%%%%%%%%%%%%%%%%%%%%%%%%%%
\subsection{Communality}
\label{subsec:factor-analysis-communality}
%%%%%%%%%%%%%%%%%%%%%%%%%%%%%%%%%%%%%%%%%%%%%%%%%%%%%%%%%%%%

For variable \(j\), the communality represents the variance explained by
the common factors.

If

\[
\lambda_{jr}
\]

is the loading of variable \(j\) on factor \(r\), then

\[
\boxed{
h_j^2
=
\sum_r
\lambda_{jr}^2
}
\]

under a standardized orthogonal-factor model.

The remaining variance is uniqueness.

%%%%%%%%%%%%%%%%%%%%%%%%%%%%%%%%%%%%%%%%%%%%%%%%%%%%%%%%%%%%
\subsection{PCA Versus Factor Analysis}
\label{subsec:pca-versus-factor-analysis}
%%%%%%%%%%%%%%%%%%%%%%%%%%%%%%%%%%%%%%%%%%%%%%%%%%%%%%%%%%%%

PCA decomposes total observed variance.

Factor analysis posits a latent-variable model separating common and
unique variation.

Thus:

\[
\boxed{
\text{PCA}
=
\text{variance decomposition},
}
\]

whereas

\[
\boxed{
\text{factor analysis}
=
\text{latent covariance model}.
}
\]

They are related but not equivalent.

%%%%%%%%%%%%%%%%%%%%%%%%%%%%%%%%%%%%%%%%%%%%%%%%%%%%%%%%%%%%
\subsection{Factor Rotation}
\label{subsec:factor-rotation}
%%%%%%%%%%%%%%%%%%%%%%%%%%%%%%%%%%%%%%%%%%%%%%%%%%%%%%%%%%%%

Factor loadings may be rotated to improve interpretability.

An orthogonal rotation preserves perpendicular factor axes.

An oblique rotation allows factors to correlate.

Common rotation criteria attempt to produce patterns where each variable
loads strongly on only a few factors.

%%%%%%%%%%%%%%%%%%%%%%%%%%%%%%%%%%%%%%%%%%%%%%%%%%%%%%%%%%%%
\subsection{Linear Discriminant Analysis}
\label{subsec:linear-discriminant-analysis}
%%%%%%%%%%%%%%%%%%%%%%%%%%%%%%%%%%%%%%%%%%%%%%%%%%%%%%%%%%%%

Linear discriminant analysis, abbreviated LDA, is a classification method
for multivariate predictors.

Suppose class \(k\) follows

\[
\boxed{
\mathbf{X}
\mid
Y=k
\sim
N_p
(
\boldsymbol{\mu}_k,
\Sigma
),
}
\]

with a covariance matrix shared by all classes.

Bayes' rule then produces linear decision boundaries.

%%%%%%%%%%%%%%%%%%%%%%%%%%%%%%%%%%%%%%%%%%%%%%%%%%%%%%%%%%%%
\subsection{LDA Discriminant Function}
\label{subsec:lda-discriminant-function}
%%%%%%%%%%%%%%%%%%%%%%%%%%%%%%%%%%%%%%%%%%%%%%%%%%%%%%%%%%%%

For class \(k\) with prior probability \(\pi_k\), the LDA discriminant
function is

\[
\boxed{
\delta_k(\mathbf{x})
=
\mathbf{x}^T
\Sigma^{-1}
\boldsymbol{\mu}_k
-
\frac12
\boldsymbol{\mu}_k^T
\Sigma^{-1}
\boldsymbol{\mu}_k
+
\log\pi_k.
}
\]

Assign \(\mathbf{x}\) to the class having the largest
\(\delta_k(\mathbf{x})\).

%%%%%%%%%%%%%%%%%%%%%%%%%%%%%%%%%%%%%%%%%%%%%%%%%%%%%%%%%%%%
\subsection{Why LDA Boundaries Are Linear}
\label{subsec:why-lda-linear}
%%%%%%%%%%%%%%%%%%%%%%%%%%%%%%%%%%%%%%%%%%%%%%%%%%%%%%%%%%%%

When comparing two class discriminant functions, the common quadratic
term involving

\[
\mathbf{x}^T\Sigma^{-1}\mathbf{x}
\]

cancels.

The resulting class boundary is therefore linear in \(\mathbf{x}\).

This is the origin of the name linear discriminant analysis.

%%%%%%%%%%%%%%%%%%%%%%%%%%%%%%%%%%%%%%%%%%%%%%%%%%%%%%%%%%%%
\subsection{Quadratic Discriminant Analysis}
\label{subsec:quadratic-discriminant-analysis}
%%%%%%%%%%%%%%%%%%%%%%%%%%%%%%%%%%%%%%%%%%%%%%%%%%%%%%%%%%%%

Quadratic discriminant analysis, abbreviated QDA, allows each class to
have its own covariance matrix:

\[
\boxed{
\mathbf{X}
\mid
Y=k
\sim
N_p
(
\boldsymbol{\mu}_k,
\Sigma_k
).
}
\]

Because the covariance matrices differ, quadratic terms no longer cancel.

Decision boundaries are therefore generally quadratic.

%%%%%%%%%%%%%%%%%%%%%%%%%%%%%%%%%%%%%%%%%%%%%%%%%%%%%%%%%%%%
\subsection{LDA Versus QDA}
\label{subsec:lda-versus-qda}
%%%%%%%%%%%%%%%%%%%%%%%%%%%%%%%%%%%%%%%%%%%%%%%%%%%%%%%%%%%%

LDA estimates one common covariance matrix.

QDA estimates a separate covariance matrix for every class.

Therefore QDA is more flexible but requires more data.

This creates a bias--variance tradeoff:

\[
\boxed{
\text{LDA}
=
\text{more structure, lower variance},
}
\]

\[
\boxed{
\text{QDA}
=
\text{less structure, greater flexibility}.
}
\]

%%%%%%%%%%%%%%%%%%%%%%%%%%%%%%%%%%%%%%%%%%%%%%%%%%%%%%%%%%%%
\subsection{Fisher's Linear Discriminant}
\label{subsec:fisher-linear-discriminant}
%%%%%%%%%%%%%%%%%%%%%%%%%%%%%%%%%%%%%%%%%%%%%%%%%%%%%%%%%%%%

For two groups, Fisher's approach searches for a direction
\(\mathbf{a}\) maximizing separation of projected group means relative to
within-group variability.

A criterion is

\[
\boxed{
J(\mathbf{a})
=
\frac{
[
\mathbf{a}^T
(
\boldsymbol{\mu}_1-\boldsymbol{\mu}_2
)
]^2
}{
\mathbf{a}^T
S_W
\mathbf{a}
}.
}
\]

The maximizing direction is proportional to

\[
\boxed{
S_W^{-1}
(
\boldsymbol{\mu}_1-\boldsymbol{\mu}_2
).
}
\]

%%%%%%%%%%%%%%%%%%%%%%%%%%%%%%%%%%%%%%%%%%%%%%%%%%%%%%%%%%%%
\subsection{Canonical Correlation Analysis}
\label{subsec:canonical-correlation-analysis}
%%%%%%%%%%%%%%%%%%%%%%%%%%%%%%%%%%%%%%%%%%%%%%%%%%%%%%%%%%%%

Canonical correlation analysis studies association between two sets of
variables.

Suppose

\[
\mathbf{X}\in\mathbb{R}^p
\]

and

\[
\mathbf{Y}\in\mathbb{R}^q.
\]

Choose linear combinations

\[
\boxed{
U
=
\mathbf{a}^T\mathbf{X},
}
\]

and

\[
\boxed{
V
=
\mathbf{b}^T\mathbf{Y}
}
\]

to maximize

\[
\boxed{
\operatorname{Corr}(U,V).
}
\]

%%%%%%%%%%%%%%%%%%%%%%%%%%%%%%%%%%%%%%%%%%%%%%%%%%%%%%%%%%%%
\subsection{Canonical Variates}
\label{subsec:canonical-variates}
%%%%%%%%%%%%%%%%%%%%%%%%%%%%%%%%%%%%%%%%%%%%%%%%%%%%%%%%%%%%

The maximizing combinations

\[
U_1
\]

and

\[
V_1
\]

form the first pair of canonical variates.

Subsequent pairs maximize correlation subject to orthogonality-type
constraints with previous canonical variates.

The corresponding correlations are called canonical correlations.

%%%%%%%%%%%%%%%%%%%%%%%%%%%%%%%%%%%%%%%%%%%%%%%%%%%%%%%%%%%%
\subsection{Clustering Preview}
\label{subsec:multivariate-clustering-preview}
%%%%%%%%%%%%%%%%%%%%%%%%%%%%%%%%%%%%%%%%%%%%%%%%%%%%%%%%%%%%

Cluster analysis attempts to divide observations into groups using
similarity in multivariate space.

Unlike LDA, clustering does not begin with known class labels.

Thus it is an

\[
\boxed{
\text{unsupervised learning}
}
\]

problem.

Common methods include:

\[
\boxed{
k\text{-means}
}
\]

and

\[
\boxed{
\text{hierarchical clustering}.
}
\]

These are developed more fully in Statistical Learning.

%%%%%%%%%%%%%%%%%%%%%%%%%%%%%%%%%%%%%%%%%%%%%%%%%%%%%%%%%%%%
\subsection{\(k\)-Means Objective Preview}
\label{subsec:kmeans-preview}
%%%%%%%%%%%%%%%%%%%%%%%%%%%%%%%%%%%%%%%%%%%%%%%%%%%%%%%%%%%%

Suppose observations are divided into clusters

\[
C_1,\ldots,C_K.
\]

The \(k\)-means objective minimizes within-cluster squared Euclidean
distance:

\[
\boxed{
\sum_{k=1}^{K}
\sum_{\mathbf{x}_i\in C_k}
\|
\mathbf{x}_i
-
\boldsymbol{\mu}_k
\|_2^2.
}
\]

The cluster center

\[
\boldsymbol{\mu}_k
\]

is the mean vector of cluster \(k\).

%%%%%%%%%%%%%%%%%%%%%%%%%%%%%%%%%%%%%%%%%%%%%%%%%%%%%%%%%%%%
\subsection{Scaling in Multivariate Analysis}
\label{subsec:multivariate-scaling}
%%%%%%%%%%%%%%%%%%%%%%%%%%%%%%%%%%%%%%%%%%%%%%%%%%%%%%%%%%%%

Suppose one variable is measured in dollars with standard deviation

\[
10000
\]

while another is measured on a scale with standard deviation

\[
2.
\]

Distance- and variance-based procedures can become dominated by the
first variable.

Thus standardization may be essential before applying:

\[
\boxed{
\text{PCA},
}
\]

\[
\boxed{
\text{clustering},
}
\]

or

\[
\boxed{
\text{distance-based methods}.
}
\]

%%%%%%%%%%%%%%%%%%%%%%%%%%%%%%%%%%%%%%%%%%%%%%%%%%%%%%%%%%%%
\subsection{Multivariate Standardization}
\label{subsec:multivariate-standardization}
%%%%%%%%%%%%%%%%%%%%%%%%%%%%%%%%%%%%%%%%%%%%%%%%%%%%%%%%%%%%

Componentwise standardization uses

\[
\boxed{
Z_j
=
\frac{
X_j-\mu_j
}{
\sigma_j
}.
}
\]

After standardization,

\[
\boxed{
\mathbb{E}[Z_j]=0,
}
\]

and

\[
\boxed{
\operatorname{Var}(Z_j)=1.
}
\]

The covariance matrix of the standardized vector is the correlation
matrix.

%%%%%%%%%%%%%%%%%%%%%%%%%%%%%%%%%%%%%%%%%%%%%%%%%%%%%%%%%%%%
\subsection{Whitening}
\label{subsec:multivariate-whitening}
%%%%%%%%%%%%%%%%%%%%%%%%%%%%%%%%%%%%%%%%%%%%%%%%%%%%%%%%%%%%

Whitening transforms a vector so its covariance becomes the identity.

If

\[
\Sigma
=
Q\Lambda Q^T
\]

is an eigenvalue decomposition with positive eigenvalues, define

\[
\boxed{
\mathbf{Z}
=
\Lambda^{-1/2}
Q^T
(
\mathbf{X}-\boldsymbol{\mu}
).
}
\]

Then

\[
\boxed{
\operatorname{Cov}(\mathbf{Z})=I.
}
\]

%%%%%%%%%%%%%%%%%%%%%%%%%%%%%%%%%%%%%%%%%%%%%%%%%%%%%%%%%%%%
\subsection{Covariance Matrix Eigen-Decomposition}
\label{subsec:covariance-eigendecomposition}
%%%%%%%%%%%%%%%%%%%%%%%%%%%%%%%%%%%%%%%%%%%%%%%%%%%%%%%%%%%%

Because \(\Sigma\) is symmetric,

\[
\boxed{
\Sigma
=
Q\Lambda Q^T,
}
\]

where

\[
Q
\]

contains orthonormal eigenvectors and

\[
\Lambda
\]

is diagonal:

\[
\boxed{
\Lambda
=
\operatorname{diag}
(
\lambda_1,\ldots,\lambda_p
).
}
\]

This decomposition underlies PCA, whitening, and multivariate normal
geometry.

%%%%%%%%%%%%%%%%%%%%%%%%%%%%%%%%%%%%%%%%%%%%%%%%%%%%%%%%%%%%
\subsection{Wishart Distribution Preview}
\label{subsec:wishart-preview}
%%%%%%%%%%%%%%%%%%%%%%%%%%%%%%%%%%%%%%%%%%%%%%%%%%%%%%%%%%%%

For multivariate normal samples, the sample covariance matrix has a
matrix-valued analogue of the chi-square distribution.

Specifically,

\[
\boxed{
(n-1)S
}
\]

has a Wishart distribution with scale matrix \(\Sigma\) and

\[
n-1
\]

degrees of freedom under the standard parameterization.

The Wishart distribution is fundamental to classical multivariate
inference.

%%%%%%%%%%%%%%%%%%%%%%%%%%%%%%%%%%%%%%%%%%%%%%%%%%%%%%%%%%%%
\subsection{Connection to the Chi-Square Distribution}
\label{subsec:wishart-chi-square-connection}
%%%%%%%%%%%%%%%%%%%%%%%%%%%%%%%%%%%%%%%%%%%%%%%%%%%%%%%%%%%%

When

\[
p=1,
\]

the covariance matrix reduces to a scalar variance.

The Wishart result reduces to

\[
\boxed{
\frac{
(n-1)S^2
}{
\sigma^2
}
\sim
\chi^2_{n-1}.
}
\]

Thus the Wishart distribution generalizes the chi-square variance result
to covariance matrices.

%%%%%%%%%%%%%%%%%%%%%%%%%%%%%%%%%%%%%%%%%%%%%%%%%%%%%%%%%%%%
\subsection{Testing Covariance Matrices}
\label{subsec:testing-covariance-matrices}
%%%%%%%%%%%%%%%%%%%%%%%%%%%%%%%%%%%%%%%%%%%%%%%%%%%%%%%%%%%%

Multivariate problems may involve hypotheses such as

\[
\boxed{
H_0:
\Sigma=\Sigma_0
}
\]

or

\[
\boxed{
H_0:
\Sigma_1=\cdots=\Sigma_g.
}
\]

Likelihood-ratio procedures can be developed under multivariate
normality.

Such tests can be highly sensitive to nonnormality.

%%%%%%%%%%%%%%%%%%%%%%%%%%%%%%%%%%%%%%%%%%%%%%%%%%%%%%%%%%%%
\subsection{Box's \(M\) Test Preview}
\label{subsec:box-m-test-preview}
%%%%%%%%%%%%%%%%%%%%%%%%%%%%%%%%%%%%%%%%%%%%%%%%%%%%%%%%%%%%

Box's \(M\) test is a classical procedure for testing equality of
covariance matrices across groups.

It is often used in discussions of MANOVA and discriminant analysis.

However, the procedure can be sensitive to deviations from multivariate
normality.

Therefore practical analysis should not rely mechanically on this test
alone.

%%%%%%%%%%%%%%%%%%%%%%%%%%%%%%%%%%%%%%%%%%%%%%%%%%%%%%%%%%%%
\subsection{High-Dimensional Covariance Problems}
\label{subsec:high-dimensional-covariance}
%%%%%%%%%%%%%%%%%%%%%%%%%%%%%%%%%%%%%%%%%%%%%%%%%%%%%%%%%%%%

If

\[
p\geq n,
\]

the ordinary sample covariance matrix cannot have full rank.

Indeed,

\[
\boxed{
\operatorname{rank}(S)
\leq
n-1.
}
\]

Therefore,

\[
\boxed{
S^{-1}
}
\]

does not exist when

\[
p\geq n
\]

under ordinary centered sampling.

This causes classical multivariate procedures to fail.

%%%%%%%%%%%%%%%%%%%%%%%%%%%%%%%%%%%%%%%%%%%%%%%%%%%%%%%%%%%%
\subsection{Curse of Dimensionality}
\label{subsec:curse-dimensionality-multivariate}
%%%%%%%%%%%%%%%%%%%%%%%%%%%%%%%%%%%%%%%%%%%%%%%%%%%%%%%%%%%%

As dimension \(p\) grows, the amount of data needed to estimate
multivariate structure grows rapidly.

Consequences include:

\[
\boxed{
\text{unstable covariance estimates},
}
\]

\[
\boxed{
\text{noisy distances},
}
\]

\[
\boxed{
\text{overfitting},
}
\]

and

\[
\boxed{
\text{increased computational cost}.
}
\]

This phenomenon is broadly called the curse of dimensionality.

%%%%%%%%%%%%%%%%%%%%%%%%%%%%%%%%%%%%%%%%%%%%%%%%%%%%%%%%%%%%
\subsection{Regularized Covariance Estimation}
\label{subsec:regularized-covariance-estimation}
%%%%%%%%%%%%%%%%%%%%%%%%%%%%%%%%%%%%%%%%%%%%%%%%%%%%%%%%%%%%

A regularized covariance estimator shrinks the sample covariance toward a
simpler matrix.

One generic form is

\[
\boxed{
\widehat{\Sigma}_{\lambda}
=
(1-\lambda)S
+
\lambda T,
}
\]

where

\[
T
\]

is a target matrix.

A common target is a multiple of the identity.

Regularization can improve stability in high dimensions.

%%%%%%%%%%%%%%%%%%%%%%%%%%%%%%%%%%%%%%%%%%%%%%%%%%%%%%%%%%%%
\subsection{Robust Multivariate Statistics Preview}
\label{subsec:robust-multivariate-preview}
%%%%%%%%%%%%%%%%%%%%%%%%%%%%%%%%%%%%%%%%%%%%%%%%%%%%%%%%%%%%

The ordinary sample mean and covariance matrix can be sensitive to
multivariate outliers.

Robust multivariate methods replace them with more resistant estimators.

Examples include robust location estimates and robust covariance
estimators.

These methods are especially useful because one unusual observation can
distort many covariance entries simultaneously.

%%%%%%%%%%%%%%%%%%%%%%%%%%%%%%%%%%%%%%%%%%%%%%%%%%%%%%%%%%%%
\subsection{Missing Multivariate Data}
\label{subsec:missing-multivariate-data}
%%%%%%%%%%%%%%%%%%%%%%%%%%%%%%%%%%%%%%%%%%%%%%%%%%%%%%%%%%%%

With \(p\) variables, missingness can occur in many patterns.

One row may be complete while another is missing only variable \(X_3\).

Deleting every row containing any missing value is called complete-case
analysis.

This can waste substantial information and may introduce bias unless the
missingness mechanism justifies it.

%%%%%%%%%%%%%%%%%%%%%%%%%%%%%%%%%%%%%%%%%%%%%%%%%%%%%%%%%%%%
\subsection{Pairwise Deletion}
\label{subsec:pairwise-deletion}
%%%%%%%%%%%%%%%%%%%%%%%%%%%%%%%%%%%%%%%%%%%%%%%%%%%%%%%%%%%%

Pairwise deletion computes each covariance using all observations
available for that pair of variables.

Although this uses more observations, the resulting covariance matrix may
have undesirable properties and may even fail to be positive
semidefinite.

Thus pairwise deletion should be used cautiously.

%%%%%%%%%%%%%%%%%%%%%%%%%%%%%%%%%%%%%%%%%%%%%%%%%%%%%%%%%%%%
\subsection{Imputation Preview}
\label{subsec:multivariate-imputation-preview}
%%%%%%%%%%%%%%%%%%%%%%%%%%%%%%%%%%%%%%%%%%%%%%%%%%%%%%%%%%%%

Missing values may be filled using an imputation model.

Single imputation treats one completed value as known and can
underrepresent uncertainty.

Multiple imputation generates several completed datasets and combines the
resulting analyses.

More complete missing-data theory is developed in computational and
mathematical statistics.

%%%%%%%%%%%%%%%%%%%%%%%%%%%%%%%%%%%%%%%%%%%%%%%%%%%%%%%%%%%%
\subsection{Visualization of Multivariate Data}
\label{subsec:multivariate-visualization}
%%%%%%%%%%%%%%%%%%%%%%%%%%%%%%%%%%%%%%%%%%%%%%%%%%%%%%%%%%%%

Useful visualizations include:

\[
\boxed{
\text{scatterplot matrices},
}
\]

\[
\boxed{
\text{correlation matrices},
}
\]

\[
\boxed{
\text{parallel-coordinate plots},
}
\]

\[
\boxed{
\text{PCA score plots},
}
\]

and

\[
\boxed{
\text{biplots}.
}
\]

No single visualization captures all high-dimensional structure.

%%%%%%%%%%%%%%%%%%%%%%%%%%%%%%%%%%%%%%%%%%%%%%%%%%%%%%%%%%%%
\subsection{Scatterplot Matrix}
\label{subsec:scatterplot-matrix}
%%%%%%%%%%%%%%%%%%%%%%%%%%%%%%%%%%%%%%%%%%%%%%%%%%%%%%%%%%%%

A scatterplot matrix displays pairwise scatterplots for several
variables.

It can reveal:

\[
\boxed{
\text{linear relationships},
}
\]

\[
\boxed{
\text{nonlinearity},
}
\]

\[
\boxed{
\text{clusters},
}
\]

and

\[
\boxed{
\text{outliers}.
}
\]

However, purely higher-dimensional relationships may remain hidden.

%%%%%%%%%%%%%%%%%%%%%%%%%%%%%%%%%%%%%%%%%%%%%%%%%%%%%%%%%%%%
\subsection{Biplot Preview}
\label{subsec:pca-biplot}
%%%%%%%%%%%%%%%%%%%%%%%%%%%%%%%%%%%%%%%%%%%%%%%%%%%%%%%%%%%%

A PCA biplot displays observations and variable directions in a reduced
principal-component space.

Observation points represent component scores.

Variable vectors represent loading information.

Interpretation depends on the specific scaling convention used to
construct the biplot.

%%%%%%%%%%%%%%%%%%%%%%%%%%%%%%%%%%%%%%%%%%%%%%%%%%%%%%%%%%%%
\subsection{Multiple Testing Across Variables}
\label{subsec:multivariate-multiple-testing}
%%%%%%%%%%%%%%%%%%%%%%%%%%%%%%%%%%%%%%%%%%%%%%%%%%%%%%%%%%%%

If \(p\) variables are tested separately, the probability of false
positive findings increases.

Possible responses include:

\[
\boxed{
\text{Bonferroni correction},
}
\]

\[
\boxed{
\text{false-discovery-rate control},
}
\]

or

\[
\boxed{
\text{joint multivariate tests}.
}
\]

The correct approach depends on the scientific goal.

%%%%%%%%%%%%%%%%%%%%%%%%%%%%%%%%%%%%%%%%%%%%%%%%%%%%%%%%%%%%
\subsection{Multivariate Effect Interpretation}
\label{subsec:multivariate-effect-interpretation}
%%%%%%%%%%%%%%%%%%%%%%%%%%%%%%%%%%%%%%%%%%%%%%%%%%%%%%%%%%%%

A multivariate test may indicate that groups differ in the vector of
means.

This does not identify which variables are responsible.

Follow-up analysis may include:

\[
\boxed{
\text{individual confidence intervals},
}
\]

\[
\boxed{
\text{contrasts},
}
\]

\[
\boxed{
\text{discriminant directions},
}
\]

or

\[
\boxed{
\text{dimension-reduction methods}.
}
\]

These follow-up analyses should respect multiplicity and study design.

%%%%%%%%%%%%%%%%%%%%%%%%%%%%%%%%%%%%%%%%%%%%%%%%%%%%%%%%%%%%
\subsection{Multivariate Statistics Workflow}
\label{subsec:multivariate-statistics-workflow}
%%%%%%%%%%%%%%%%%%%%%%%%%%%%%%%%%%%%%%%%%%%%%%%%%%%%%%%%%%%%

A useful workflow is:

\begin{enumerate}
    \item identify all variables and their units;
    \item determine the multivariate scientific question;
    \item inspect missing values;
    \item examine univariate distributions;
    \item inspect pairwise scatterplots;
    \item compute covariance and correlation matrices;
    \item identify scale differences;
    \item standardize variables when appropriate;
    \item inspect multivariate outliers;
    \item assess sample size relative to dimension;
    \item select an inferential or dimension-reduction method;
    \item verify covariance and distributional assumptions;
    \item fit the multivariate model;
    \item examine joint tests before extensive follow-up testing;
    \item interpret multivariate directions or components carefully;
    \item validate conclusions using independent or held-out data when
          prediction is involved.
\end{enumerate}

%%%%%%%%%%%%%%%%%%%%%%%%%%%%%%%%%%%%%%%%%%%%%%%%%%%%%%%%%%%%
\subsection{Choosing a Multivariate Method}
\label{subsec:choosing-multivariate-method}
%%%%%%%%%%%%%%%%%%%%%%%%%%%%%%%%%%%%%%%%%%%%%%%%%%%%%%%%%%%%

To compare one population mean vector with a known vector:

\[
\boxed{
\text{one-sample Hotelling }T^2.
}
\]

To compare two multivariate group means:

\[
\boxed{
\text{two-sample Hotelling }T^2.
}
\]

To compare several groups on several responses:

\[
\boxed{
\text{MANOVA}.
}
\]

To reduce dimension without class labels:

\[
\boxed{
\text{PCA}.
}
\]

To model latent common factors:

\[
\boxed{
\text{factor analysis}.
}
\]

To classify known groups:

\[
\boxed{
\text{LDA or QDA}.
}
\]

To study association between two variable sets:

\[
\boxed{
\text{canonical correlation}.
}
\]

%%%%%%%%%%%%%%%%%%%%%%%%%%%%%%%%%%%%%%%%%%%%%%%%%%%%%%%%%%%%
\subsection{Common Errors}
\label{subsec:multivariate-common-errors}
%%%%%%%%%%%%%%%%%%%%%%%%%%%%%%%%%%%%%%%%%%%%%%%%%%%%%%%%%%%%

\subsubsection{Analyzing Every Variable Separately Without Multiplicity Control}

Separate univariate tests can inflate the probability of false
discoveries.

Joint multivariate structure may also be lost.

%%%%%%%%%%%%%%%%%%%%%%%%%%%%%%%%%%%%%%%%%%%%%%%%%%%%%%%%%%%%
\subsubsection{Ignoring Variable Scale}

Variance- and distance-based methods can be dominated by variables
measured on large numerical scales.

Standardization may be required.

%%%%%%%%%%%%%%%%%%%%%%%%%%%%%%%%%%%%%%%%%%%%%%%%%%%%%%%%%%%%
\subsubsection{Using Correlation When Covariance Is Scientifically Relevant}

Standardizing removes unit-scale differences.

If absolute variation is scientifically meaningful, covariance-based
analysis may be preferable.

%%%%%%%%%%%%%%%%%%%%%%%%%%%%%%%%%%%%%%%%%%%%%%%%%%%%%%%%%%%%
\subsubsection{Using Covariance PCA with Incommensurate Units Without Thought}

A variable measured in thousands may dominate a variable measured in
small units even when the latter is scientifically important.

%%%%%%%%%%%%%%%%%%%%%%%%%%%%%%%%%%%%%%%%%%%%%%%%%%%%%%%%%%%%
\subsubsection{Treating PCA Components as Causal Variables}

Principal components are algebraically constructed directions of
variation.

They do not automatically represent underlying causal mechanisms.

%%%%%%%%%%%%%%%%%%%%%%%%%%%%%%%%%%%%%%%%%%%%%%%%%%%%%%%%%%%%
\subsubsection{Interpreting PCA as Maximizing Prediction}

PCA maximizes predictor variance, not predictive relevance to an external
response.

%%%%%%%%%%%%%%%%%%%%%%%%%%%%%%%%%%%%%%%%%%%%%%%%%%%%%%%%%%%%
\subsubsection{Treating PCA and Factor Analysis as Identical}

PCA decomposes total variation.

Factor analysis models covariance through latent common factors and
unique variation.

%%%%%%%%%%%%%%%%%%%%%%%%%%%%%%%%%%%%%%%%%%%%%%%%%%%%%%%%%%%%
\subsubsection{Ignoring Multivariate Outliers}

An observation may be unusual jointly even when none of its individual
coordinates appears extreme.

%%%%%%%%%%%%%%%%%%%%%%%%%%%%%%%%%%%%%%%%%%%%%%%%%%%%%%%%%%%%
\subsubsection{Using Mahalanobis Distance with an Unstable Covariance Estimate}

If \(S\) is poorly estimated, then

\[
S^{-1}
\]

can be extremely unstable.

This is especially dangerous when \(p\) is large relative to \(n\).

%%%%%%%%%%%%%%%%%%%%%%%%%%%%%%%%%%%%%%%%%%%%%%%%%%%%%%%%%%%%
\subsubsection{Attempting to Invert a Singular Covariance Matrix}

When

\[
p\geq n,
\]

the ordinary centered sample covariance matrix is singular.

Classical inverse-covariance methods require modification or
regularization.

%%%%%%%%%%%%%%%%%%%%%%%%%%%%%%%%%%%%%%%%%%%%%%%%%%%%%%%%%%%%
\subsubsection{Assuming Zero Correlation Means Independence}

Zero covariance or correlation does not generally imply independence.

The implication holds under special models such as joint multivariate
normality.

%%%%%%%%%%%%%%%%%%%%%%%%%%%%%%%%%%%%%%%%%%%%%%%%%%%%%%%%%%%%
\subsubsection{Ignoring Covariance Differences in Classification}

LDA assumes a common within-class covariance matrix.

If class covariance structures differ substantially, QDA or another
method may be more appropriate.

%%%%%%%%%%%%%%%%%%%%%%%%%%%%%%%%%%%%%%%%%%%%%%%%%%%%%%%%%%%%
\subsubsection{Using QDA with Too Little Data}

QDA estimates a separate covariance matrix for each class.

This requires considerably more data than LDA.

%%%%%%%%%%%%%%%%%%%%%%%%%%%%%%%%%%%%%%%%%%%%%%%%%%%%%%%%%%%%
\subsubsection{Interpreting a Significant MANOVA Without Follow-Up Analysis}

A significant MANOVA only establishes a multivariate difference.

It does not identify the specific responses or groups responsible.

%%%%%%%%%%%%%%%%%%%%%%%%%%%%%%%%%%%%%%%%%%%%%%%%%%%%%%%%%%%%
\subsubsection{Ignoring Dependence or Repeated Measures}

Classical multivariate methods still require the observational dependence
structure to be modeled correctly.

Multiple variables on one subject do not create independent observations.

%%%%%%%%%%%%%%%%%%%%%%%%%%%%%%%%%%%%%%%%%%%%%%%%%%%%%%%%%%%%
\subsection{Applications}
\label{subsec:multivariate-applications}
%%%%%%%%%%%%%%%%%%%%%%%%%%%%%%%%%%%%%%%%%%%%%%%%%%%%%%%%%%%%

\begin{applicationbox}

\textbf{Environmental Science.}
Temperature, dissolved oxygen, depth, salinity, nutrient levels, and time
may be analyzed jointly to understand environmental systems.

\medskip

\textbf{Medicine.}
Multiple biomarkers, physiological measurements, and clinical outcomes
often form correlated response vectors.

\medskip

\textbf{Education Research.}
Student performance may be summarized through several grades, assessment
scores, attendance measures, and survey variables.

\medskip

\textbf{Finance.}
Covariance matrices describe relationships among asset returns and are
central to portfolio risk.

\medskip

\textbf{Engineering.}
Multivariate process-monitoring systems track several correlated quality
measurements simultaneously.

\medskip

\textbf{Psychology.}
Factor analysis and PCA are used to study latent structures underlying
multiple observed measurements.

\medskip

\textbf{Biology.}
Multivariate methods analyze gene expression, morphology, ecological
traits, and high-dimensional biological measurements.

\medskip

\textbf{Machine Learning.}
PCA, discriminant analysis, covariance regularization, and clustering are
fundamental tools for dimension reduction, classification, and
unsupervised learning.

\end{applicationbox}

%%%%%%%%%%%%%%%%%%%%%%%%%%%%%%%%%%%%%%%%%%%%%%%%%%%%%%%%%%%%
\subsection{Exercises}
\label{subsec:multivariate-exercises}
%%%%%%%%%%%%%%%%%%%%%%%%%%%%%%%%%%%%%%%%%%%%%%%%%%%%%%%%%%%%

\begin{exercise}[1]
Define a random vector.
\end{exercise}

\begin{exercise}[1]
Define the population mean vector.
\end{exercise}

\begin{exercise}[1]
Define the covariance matrix.
\end{exercise}

\begin{exercise}[1]
Explain why a covariance matrix is symmetric.
\end{exercise}

\begin{exercise}[1]
Define a correlation matrix.
\end{exercise}

\begin{exercise}[1]
Define the sample mean vector.
\end{exercise}

\begin{exercise}[1]
Define the sample covariance matrix.
\end{exercise}

\begin{exercise}[1]
Define Mahalanobis distance.
\end{exercise}

\begin{exercise}[1]
State the one-sample Hotelling \(T^2\) statistic.
\end{exercise}

\begin{exercise}[1]
Define MANOVA.
\end{exercise}

\begin{exercise}[1]
Define a principal component.
\end{exercise}

\begin{exercise}[1]
Define explained variance in PCA.
\end{exercise}

\begin{exercise}[1]
Distinguish PCA and factor analysis.
\end{exercise}

\begin{exercise}[1]
Distinguish LDA and QDA.
\end{exercise}

\begin{exercise}[2]
For

\[
\Sigma
=
\begin{pmatrix}
4&2\\
2&9
\end{pmatrix},
\]

identify the two variances and covariance.
\end{exercise}

\begin{exercise}[2]
Compute the correlation corresponding to the covariance matrix

\[
\Sigma
=
\begin{pmatrix}
4&2\\
2&9
\end{pmatrix}.
\]
\end{exercise}

\begin{exercise}[2]
Suppose

\[
\mathbf{a}
=
\begin{pmatrix}
1\\
2
\end{pmatrix}
\]

and

\[
\Sigma
=
\begin{pmatrix}
2&1\\
1&3
\end{pmatrix}.
\]

Compute

\[
\operatorname{Var}
(
\mathbf{a}^T\mathbf{X}
).
\]
\end{exercise}

\begin{exercise}[2]
For observations

\[
(1,2),
\quad
(3,4),
\quad
(5,6),
\]

compute the sample mean vector.
\end{exercise}

\begin{exercise}[2]
For

\[
\boldsymbol{\mu}
=
\begin{pmatrix}
0\\
0
\end{pmatrix},
\qquad
\Sigma=I,
\]

compute the Mahalanobis distance of

\[
\mathbf{x}
=
\begin{pmatrix}
3\\
4
\end{pmatrix}.
\]
\end{exercise}

\begin{exercise}[2]
If PCA eigenvalues are

\[
10,\ 5,\ 3,\ 2,
\]

compute the proportion of variance explained by the first component.
\end{exercise}

\begin{exercise}[2]
For the previous eigenvalues, compute the cumulative proportion explained
by the first two components.
\end{exercise}

\begin{exercise}[2]
Suppose \(p=4\) and \(n=30\). State the numerator and denominator degrees
of freedom for the transformed one-sample Hotelling test.
\end{exercise}

\begin{exercise}[2]
Suppose \(p=12\) and \(n=10\). Explain why the ordinary sample covariance
matrix cannot be inverted.
\end{exercise}

\begin{exercise}[2]
If an LDA coefficient is based on two classes having equal covariance
matrices, explain why the decision boundary is linear.
\end{exercise}

\begin{exercise}[3]
Prove that every covariance matrix is positive semidefinite.
\end{exercise}

\begin{exercise}[3]
Show that

\[
\operatorname{Cov}(A\mathbf{X})
=
A\Sigma A^T.
\]
\end{exercise}

\begin{exercise}[3]
Derive the variance

\[
\mathbf{a}^T
\Sigma
\mathbf{a}
\]

for a linear combination.
\end{exercise}

\begin{exercise}[3]
Show that standardizing a random vector transforms its covariance matrix
into its correlation matrix.
\end{exercise}

\begin{exercise}[3]
Verify that the multivariate normal density reduces to the ordinary
normal density when \(p=1\).
\end{exercise}

\begin{exercise}[3]
Show that multivariate normal density contours are ellipsoids.
\end{exercise}

\begin{exercise}[3]
Explain why zero covariance implies independence under joint normality
but not in general.
\end{exercise}

\begin{exercise}[3]
Show that squared Mahalanobis distance is invariant under nonsingular
linear changes of coordinates.
\end{exercise}

\begin{exercise}[3]
Show that Hotelling \(T^2\) reduces to \(t^2\) when \(p=1\).
\end{exercise}

\begin{exercise}[3]
Explain why Hotelling's test accounts for correlation among response
variables.
\end{exercise}

\begin{exercise}[3]
Compare Bonferroni intervals with a Hotelling confidence ellipsoid.
\end{exercise}

\begin{exercise}[3]
Explain the difference between ANOVA and MANOVA.
\end{exercise}

\begin{exercise}[3]
Interpret small values of Wilks' lambda.
\end{exercise}

\begin{exercise}[3]
Explain conceptually how Pillai's trace differs from Roy's largest root.
\end{exercise}

\begin{exercise}[3]
Using a Lagrange multiplier, derive the eigenvalue equation

\[
\Sigma\mathbf{a}
=
\lambda\mathbf{a}
\]

for the first principal component.
\end{exercise}

\begin{exercise}[3]
Show that principal components are uncorrelated.
\end{exercise}

\begin{exercise}[3]
Prove that total variance equals

\[
\operatorname{tr}(\Sigma)
=
\sum_j\lambda_j.
\]
\end{exercise}

\begin{exercise}[3]
Explain why PCA based on the covariance matrix can be sensitive to
measurement units.
\end{exercise}

\begin{exercise}[3]
Explain why PCA is considered unsupervised.
\end{exercise}

\begin{exercise}[3]
Explain the difference between PCA scores and PCA loadings.
\end{exercise}

\begin{exercise}[3]
Derive the factor-model covariance

\[
\Sigma
=
\Lambda\Lambda^T+\Psi.
\]
\end{exercise}

\begin{exercise}[3]
Explain why factor rotation can improve interpretability.
\end{exercise}

\begin{exercise}[3]
Derive the two-class LDA discriminant boundary.
\end{exercise}

\begin{exercise}[3]
Explain why QDA needs more observations than LDA.
\end{exercise}

\begin{exercise}[3]
Explain the objective of canonical correlation analysis.
\end{exercise}

\begin{exercise}[3]
Explain why Euclidean distance can become misleading when variables have
very different scales.
\end{exercise}

\begin{exercise}[3]
Show that a centered \(n\times p\) data matrix has rank at most \(n-1\).
\end{exercise}

\begin{exercise}[3]
Explain why pairwise deletion can produce an invalid covariance matrix.
\end{exercise}

\begin{exercise}[4]
Derive the multivariate normal conditional distribution.
\end{exercise}

\begin{exercise}[4]
Study the Schur complement appearing in the conditional covariance
matrix.
\end{exercise}

\begin{exercise}[4]
Derive the exact one-sample Hotelling \(T^2\) distribution under
multivariate normality.
\end{exercise}

\begin{exercise}[4]
Derive the two-sample Hotelling \(T^2\) statistic.
\end{exercise}

\begin{exercise}[4]
Study simultaneous inference for all linear combinations of a mean
vector.
\end{exercise}

\begin{exercise}[4]
Derive Wilks' lambda from likelihood ratios in MANOVA.
\end{exercise}

\begin{exercise}[4]
Study Pillai's trace, Hotelling--Lawley trace, and Roy's largest root.
\end{exercise}

\begin{exercise}[4]
Derive PCA using the spectral theorem.
\end{exercise}

\begin{exercise}[4]
Prove the minimum reconstruction-error property of PCA.
\end{exercise}

\begin{exercise}[4]
Study PCA through singular-value decomposition.
\end{exercise}

\begin{exercise}[4]
Investigate probabilistic PCA.
\end{exercise}

\begin{exercise}[4]
Derive maximum-likelihood estimators in the Gaussian factor-analysis
model.
\end{exercise}

\begin{exercise}[4]
Study orthogonal and oblique factor rotations.
\end{exercise}

\begin{exercise}[4]
Derive LDA from Bayes classification under multivariate normal class
models.
\end{exercise}

\begin{exercise}[4]
Derive Fisher's discriminant direction.
\end{exercise}

\begin{exercise}[4]
Study canonical correlation through generalized eigenvalue problems.
\end{exercise}

\begin{exercise}[4]
Derive the Wishart distribution of the sample covariance matrix.
\end{exercise}

\begin{exercise}[4]
Investigate likelihood-ratio tests for covariance matrices.
\end{exercise}

\begin{exercise}[4]
Study Box's \(M\) test and its sensitivity to nonnormality.
\end{exercise}

\begin{exercise}[4]
Investigate shrinkage covariance estimators.
\end{exercise}

\begin{exercise}[4]
Study robust estimators of multivariate location and covariance.
\end{exercise}

\begin{exercise}[4]
Investigate high-dimensional PCA when \(p/n\) is not small.
\end{exercise}

%%%%%%%%%%%%%%%%%%%%%%%%%%%%%%%%%%%%%%%%%%%%%%%%%%%%%%%%%%%%
\subsection{Conceptual Summary}
\label{subsec:multivariate-summary}
%%%%%%%%%%%%%%%%%%%%%%%%%%%%%%%%%%%%%%%%%%%%%%%%%%%%%%%%%%%%

Multivariate statistics studies random vectors

\[
\boxed{
\mathbf{X}
=
\begin{pmatrix}
X_1\\
\vdots\\
X_p
\end{pmatrix}.
}
\]

The mean vector is

\[
\boxed{
\boldsymbol{\mu}
=
\mathbb{E}[\mathbf{X}],
}
\]

and the covariance matrix is

\[
\boxed{
\Sigma
=
\mathbb{E}
\left[
(
\mathbf{X}-\boldsymbol{\mu}
)
(
\mathbf{X}-\boldsymbol{\mu}
)^T
\right].
}
\]

For a linear transformation,

\[
\boxed{
\operatorname{Cov}
(
A\mathbf{X}
)
=
A\Sigma A^T.
}
\]

Mahalanobis distance is

\[
\boxed{
D_M^2
=
(
\mathbf{x}-\boldsymbol{\mu}
)^T
\Sigma^{-1}
(
\mathbf{x}-\boldsymbol{\mu}
).
}
\]

The one-sample Hotelling statistic is

\[
\boxed{
T^2
=
n
(
\overline{\mathbf{X}}
-
\boldsymbol{\mu}_0
)^T
S^{-1}
(
\overline{\mathbf{X}}
-
\boldsymbol{\mu}_0
).
}
\]

MANOVA extends ANOVA to multiple response variables.

PCA uses eigenvectors of the covariance or correlation matrix:

\[
\boxed{
\Sigma\mathbf{a}_j
=
\lambda_j\mathbf{a}_j.
}
\]

The proportion of variance explained by component \(j\) is

\[
\boxed{
\frac{
\lambda_j
}{
\sum_k\lambda_k
}.
}
\]

Factor analysis models latent common structure:

\[
\boxed{
\mathbf{X}
=
\boldsymbol{\mu}
+
\Lambda\mathbf{F}
+
\boldsymbol{\varepsilon}.
}
\]

LDA and QDA use multivariate probability models for classification.

Thus multivariate statistics combines

\[
\boxed{
\text{linear algebra}
+
\text{probability}
+
\text{covariance}
+
\text{hypothesis testing}
+
\text{dimension reduction}.
}
\]

%%%%%%%%%%%%%%%%%%%%%%%%%%%%%%%%%%%%%%%%%%%%%%%%%%%%%%%%%%%%
\subsection{Section Connections}
\label{subsec:multivariate-connections}
%%%%%%%%%%%%%%%%%%%%%%%%%%%%%%%%%%%%%%%%%%%%%%%%%%%%%%%%%%%%

\chapterconnection{
Multivariate Statistics extends the univariate and regression methods of
Sections~\ref{sec:statistical-estimation},
\ref{sec:hypothesis-testing},
\ref{sec:regression}, and
\ref{sec:analysis-of-variance} to vector-valued observations. Hotelling's
\(T^2\) generalizes the one-sample and two-sample \(t\)-tests, while
MANOVA generalizes ANOVA to several correlated responses. Principal
component analysis and factor analysis introduce dimension reduction and
latent structure, and discriminant analysis connects multivariate
probability with classification. The next section, Mathematical
Statistics, develops the probability, likelihood, sufficiency,
information, decision-theoretic, and asymptotic foundations underlying
the statistical procedures used throughout this chapter.
}

%%%%%%%%%%%%%%%%%%%%%%%%%%%%%%%%%%%%%%%%%%%%%%%%%%%%%%%%%%%%
% SECTION SOURCE TRACKING
%%%%%%%%%%%%%%%%%%%%%%%%%%%%%%%%%%%%%%%%%%%%%%%%%%%%%%%%%%%%

%------------------------------------------------------------
% 8.12 Multivariate Statistics
%------------------------------------------------------------
%
% Source basis:
%   - Standard multivariate statistical theory.
%   - Standard matrix-based statistical inference.
%   - Standard dimension-reduction and discrimination methods.
%
% Used for:
%   - random vectors;
%   - data matrices;
%   - mean vectors;
%   - covariance matrices;
%   - correlation matrices;
%   - positive semidefinite covariance structure;
%   - sample mean and covariance matrices;
%   - linear transformations;
%   - multivariate normal distributions;
%   - marginal and conditional normal distributions;
%   - Mahalanobis distance;
%   - multivariate outliers;
%   - one-sample Hotelling T-squared tests;
%   - confidence ellipsoids;
%   - simultaneous confidence intervals;
%   - two-sample Hotelling T-squared tests;
%   - MANOVA;
%   - Wilks lambda;
%   - Pillai trace;
%   - Hotelling--Lawley trace;
%   - Roy largest root;
%   - principal component analysis;
%   - eigenvalue interpretation;
%   - explained variance;
%   - PCA scores and loadings;
%   - standardization;
%   - PCA reconstruction;
%   - factor analysis;
%   - communalities;
%   - factor rotation;
%   - PCA versus factor analysis;
%   - linear discriminant analysis;
%   - quadratic discriminant analysis;
%   - Fisher discriminant analysis;
%   - canonical correlation;
%   - clustering preview;
%   - whitening;
%   - Wishart distribution preview;
%   - covariance-matrix testing;
%   - high-dimensional covariance problems;
%   - covariance regularization;
%   - robust multivariate statistics;
%   - multivariate missing data;
%   - visualization;
%   - multiple testing;
%   - applications and exercises.
%
% Notes:
%   - Mathematical Statistics is developed next in Section 8.13.
%   - Bayesian Statistics later develops multivariate posterior models.
%   - Statistical Learning develops PCA, classification, clustering, and
%     regularization more deeply.
%   - Computational Statistics develops numerical and resampling methods
%     for high-dimensional multivariate analysis.
%
%%%%%%%%%%%%%%%%%%%%%%%%%%%%%%%%%%%%%%%%%%%%%%%%%%%%%%%%%%%%